\documentclass[12pt,oneside]{book}

% ============================================================
% PAQUETES
% ============================================================
\usepackage[utf8]{inputenc}
\usepackage[T1]{fontenc}
\usepackage{lmodern}
\usepackage[spanish,es-noshorthands]{babel}
\usepackage{geometry}
\geometry{a4paper,margin=2.6cm}
\usepackage{amsmath,amssymb,mathtools}
\usepackage{graphicx}
\usepackage{float}
\usepackage{caption}
\usepackage{booktabs}
\usepackage{array}
\usepackage{longtable}
\usepackage{tabularx}
\usepackage{enumitem}
\usepackage{xcolor}
\usepackage{tcolorbox}
\tcbuselibrary{skins,breakable}
\usepackage{fancyhdr}
\usepackage{ragged2e}
\usepackage{titlesec}
\usepackage[unicode,bookmarksopen=true,breaklinks=true]{hyperref}

\definecolor{uba}{HTML}{7A1F2B}
\definecolor{ubaclaro}{HTML}{F4E9EA}
\definecolor{defcolor}{HTML}{1F4E79}
\definecolor{defbg}{HTML}{EAF1F8}
\definecolor{impcolor}{HTML}{8A5A00}
\definecolor{impbg}{HTML}{FBF0DD}
\definecolor{excolor}{HTML}{2E6B3E}
\definecolor{exbg}{HTML}{E9F5EC}
\definecolor{notecolor}{HTML}{555555}
\definecolor{notebg}{HTML}{F0F0F0}
\definecolor{examcolor}{HTML}{7A1F2B}
\definecolor{exambg}{HTML}{F6E8E9}

\hypersetup{
  colorlinks=true,
  linkcolor=uba,
  citecolor=uba,
  urlcolor=uba,
  pdftitle={Derecho Constitucional Profundizado -- Guía del Primer Parcial},
  pdfauthor={Isaac Edgar Camacho Ocampo}
}

\setcounter{tocdepth}{2}
\setlength{\parindent}{0pt}
\setlength{\parskip}{0.5em}

\titleformat{\chapter}[display]
  {\normalfont\huge\bfseries\color{uba}}
  {\Large\textcolor{uba}{Pregunta \thechapter}}{10pt}{\Huge}
\titleformat{\part}[display]
  {\normalfont\Huge\bfseries\color{uba}\centering}
  {\Large\textsc{Parte \thepart}}{20pt}{\huge}

\setlength{\headheight}{14pt}
\pagestyle{fancy}
\fancyhf{}
\renewcommand{\headrulewidth}{0.4pt}
\renewcommand{\footrulewidth}{0pt}
\fancyhead[L]{\small\textit{Derecho Constitucional Profundizado}}
\fancyhead[R]{\small\leftmark}
\fancyfoot[C]{\thepage}

% ---------- Cajas semánticas ----------
\newtcolorbox{defbox}[1][]{
  colback=defbg, colframe=defcolor, boxrule=0.6pt, arc=2pt,
  left=6pt, right=6pt, top=4pt, bottom=4pt, breakable,
  fonttitle=\bfseries, coltitle=defcolor,
  title={Definición#1},
}
\newtcolorbox{impbox}[1][Importante]{
  colback=impbg, colframe=impcolor, boxrule=0.6pt, arc=2pt,
  left=6pt, right=6pt, top=4pt, bottom=4pt, breakable,
  fonttitle=\bfseries, coltitle=impcolor,
  title={#1},
}
\newtcolorbox{exbox}[1][Ejemplo]{
  colback=exbg, colframe=excolor, boxrule=0.6pt, arc=2pt,
  left=6pt, right=6pt, top=4pt, bottom=4pt, breakable,
  fonttitle=\bfseries, coltitle=excolor,
  title={#1},
}
\newtcolorbox{notebox}[1][Nota]{
  colback=notebg, colframe=notecolor, boxrule=0.5pt, arc=2pt,
  left=6pt, right=6pt, top=3pt, bottom=3pt, breakable,
  fonttitle=\bfseries\itshape, coltitle=notecolor,
  title={#1},
}
\newtcolorbox{exambox}[1][Respuesta para el examen]{
  colback=exambg, colframe=examcolor, boxrule=0.9pt, arc=2pt,
  left=7pt, right=7pt, top=5pt, bottom=5pt, breakable,
  fonttitle=\bfseries, coltitle=examcolor,
  title={#1},
}

\newcommand{\vease}[1]{\textit{(Véase también: #1.)}}

\begin{document}

% ============================================================
% PORTADA
% ============================================================
\begin{titlepage}
\centering
\vspace*{0.5cm}
{\large\scshape Universidad de Buenos Aires (UBA)\par}
{\large\scshape Facultad de Derecho\par}
{\large\scshape Carrera de Abogacía\par}
\vspace{1.2cm}
{\color{uba}\rule{0.6\textwidth}{1pt}}\par
\vspace{0.6cm}
{\Huge\bfseries\color{uba} Derecho Constitucional\\ Profundizado\par}
\vspace{0.4cm}
{\Large Guía integral de respuestas para el\\ Primer Parcial --- Comisión 603\par}
\vspace{0.3cm}
{\large 28 de septiembre de 2026\par}
\vspace{0.6cm}
{\color{uba}\rule{0.6\textwidth}{1pt}}\par
\vspace{1cm}

{\normalsize
Interpretación constitucional y precedente judicial \\
Control de constitucionalidad, difuso y contramayoritario \\
El recurso extraordinario federal: requisitos, queja, cautelar y \textit{per saltum} \\
El principio de reserva, la objeción de conciencia, la libertad religiosa \\
y la libertad de expresión
\par}

\vspace{1.3cm}
{\itshape\large ``Comprender es el primer paso para convertir\\ el estudio en conocimiento.''\par}
\vspace{1.5cm}

{\large Docentes: Dr. Andrés Uslenghi (Adjunto) --- Dr. Darío M. Luna (JTP)\par}
\vspace{0.3cm}
{\large Apunte elaborado por: \textbf{Isaac Edgar Camacho Ocampo}\par}

\vfill
\begin{flushleft}
\small\itshape
Este es un modesto aporte a los estudiantes de ``Derecho Constitucional Profundizado''.
Reorganiza y desarrolla, en función de la guía de preguntas del primer parcial, los
apuntes integrados de la materia, contrastados con el cronograma de la cátedra y
completados —donde hacía falta— con fuentes académicas y jurisprudenciales verificables.
No pretende sustituir el material oficial de la cátedra ni la lectura de los fallos y
normas en su fuente original.
\end{flushleft}
{\footnotesize 2026\par}
\end{titlepage}

\tableofcontents

\chapter*{Presentación}
\addcontentsline{toc}{chapter}{Presentación}

Este volumen no es un apunte general de la materia, sino un desarrollo dirigido,
punto por punto, a la \textbf{guía de preguntas del primer parcial} de Derecho
Constitucional Profundizado (Comisión 603), a rendirse el 28 de septiembre de 2026.
Toma como fuente principal los apuntes integrados de la cátedra —que cubren con
solvencia la mayor parte de los catorce puntos de la guía— y los reorganiza en el
orden exacto en que la guía los presenta, profundizando cada concepto desde su base,
para que cada respuesta pueda estudiarse de manera autónoma.

En cuatro puntos de la guía los apuntes originales no traían desarrollo propio
—la distinción entre \textit{holding} y \textit{obiter dictum}, el carácter
contramayoritario del Poder Judicial, la objeción de conciencia (con los casos
``Portillo'' y ``Bahamondez'') y el derecho de réplica (caso ``Ekmekdjian
c/ Sofovich'')—. Para esos puntos se investigó en fuentes académicas y
jurisprudenciales confiables, que se citan en la bibliografía final, y se indica
expresamente que se trata de una ampliación externa y no de contenido tomado de
los apuntes de clase.

\begin{notebox}[Advertencia sobre las fuentes]
Este material combina tres fuentes: (1) los apuntes integrados de la cátedra,
que constituyen la base y la orientación de cada respuesta; (2) el cronograma de
clases de la Comisión 603, utilizado para verificar qué bibliografía y qué fallos
corresponden a cada tema; y (3) fuentes académicas externas, utilizadas
únicamente para completar los puntos de la guía que los apuntes no desarrollaban
todavía. Cuando se recurrió a una fuente externa, se lo indica expresamente en el
cuerpo del texto. No se ha inventado ningún autor, cita, fallo ni referencia
normativa: donde no fue posible verificar un dato con las fuentes disponibles,
se señala expresamente esa limitación en lugar de presentarlo como un hecho.
Este apunte no reemplaza la lectura de los fallos, normas y manuales de la
bibliografía oficial de la cátedra.
\end{notebox}

\section*{Mapa de la guía}
\addcontentsline{toc}{section}{Mapa de la guía}

La guía de catorce puntos se agrupa aquí en cuatro partes, que siguen la misma
progresión lógica que el programa de la materia:

\begin{itemize}[leftmargin=1.4em]
  \item \textbf{Parte I --- Interpretación y precedente.} Textualismo y no
  textualismo; \textit{holding} y \textit{obiter dictum} (Pregunta 1).
  \item \textbf{Parte II --- Control de constitucionalidad.} Concepto y alcance;
  supremacía constitucional (Pregunta 2); control difuso y concentrado
  (Pregunta 3); el Poder Judicial como poder contramayoritario (Pregunta 4).
  \item \textbf{Parte III --- El recurso extraordinario federal.} Cuestión
  federal simple y compleja (Pregunta 5); arbitrariedad de sentencia
  (Pregunta 6); control de oficio (Pregunta 7); requisitos del recurso
  extraordinario (Pregunta 8); recurso de queja (Pregunta 9); medida cautelar
  (Pregunta 10); \textit{per saltum} (Pregunta 11).
  \item \textbf{Parte IV --- El principio de reserva y los derechos
  individuales.} Principio de reserva y objeción de conciencia (Pregunta 12);
  libertad religiosa (Pregunta 13); libertad de expresión, censura previa,
  doctrina Campillay, real malicia y derecho de réplica (Pregunta 14).
\end{itemize}

Dentro de cada parte, cada pregunta se desarrolla siguiendo, en la medida en que
resulte útil, la progresión: contexto $\rightarrow$ concepto $\rightarrow$
fundamento $\rightarrow$ desarrollo $\rightarrow$ relación con otros temas
$\rightarrow$ jurisprudencia y ejemplos $\rightarrow$ respuesta para el examen.
Al cierre de cada parte se incluye un cuadro Cornell de repaso.


% ============================================================
% PARTE I - INTERPRETACIÓN Y PRECEDENTE
% ============================================================
\part{Interpretación constitucional y precedente judicial}

\chapter{Textualismo y no textualismo. \textit{Holding} y \textit{obiter dictum} en el precedente judicial}
\label{cap:p1}

\section{Por qué interpretar la Constitución es un problema}

Toda norma necesita ser interpretada antes de ser aplicada, pero la Constitución
plantea un problema adicional: es un texto fundacional que, a diferencia de una ley
ordinaria, no se actualiza automáticamente con el paso del tiempo. Una ley puede
reformarse con relativa facilidad cuando deja de responder a las necesidades
sociales; la Constitución, en cambio, está pensada para perdurar, y solo puede
modificarse mediante un procedimiento de reforma agravado. El resultado es una
tensión permanente entre un texto que se mantiene estable y una realidad social,
tecnológica y económica que cambia constantemente. Basta pensar que hace apenas
diez años la inteligencia artificial se percibía como ciencia ficción, y hoy plantea
problemas interpretativos concretos que ningún constituyente pudo prever.

Frente a ese problema, la primera fuente de interpretación constitucional es
siempre la letra de la norma. Pero cuando la letra por sí sola no permite resolver
el caso —porque es ambigua, porque colisiona con otra norma, o porque describe una
situación que el constituyente no imaginó— es necesario recurrir a un método, a una
\textit{corriente de interpretación}, que oriente al intérprete sobre qué peso darle
al texto, a la intención histórica de quien lo redactó, o a los fines que persigue.

\section{Tres corrientes, ordenadas según su apego a la letra}

La cátedra distingue tres grandes corrientes de interpretación constitucional,
ordenadas de mayor a menor apego literal al texto: el \textbf{textualismo}, el
\textbf{originalismo} y el \textbf{pragmatismo} (o interpretación dinámica). La
guía de examen agrupa estas corrientes bajo la dicotomía clásica
\textbf{textualismo y no textualismo}: el textualismo por un lado, y por el otro
—como variantes de una misma actitud de apartarse de la literalidad estricta— el
originalismo y el pragmatismo, cada uno por una vía distinta.

\subsection{El textualismo}

\begin{defbox}[: Textualismo]
Corriente que afianza el valor de la letra como primera fuente de interpretación y
le da un sentido estricto: la letra no es solo la primera referencia a la hora de
interpretar la norma, sino que además reduce fuertemente la creatividad del juez o
del órgano que la aplica, porque es el propio texto el que marca los límites de lo
interpretable.
\end{defbox}

En la jurisprudencia argentina, la Corte Suprema sostiene que la primera fuente de
interpretación es la letra de la norma, entendida según el uso o significado
corriente que las palabras tienen en los tiempos actuales. Se trata de una
modulación más flexible que el textualismo puro norteamericano: la Corte argentina
no congela el significado de las palabras en el momento histórico de la sanción,
sino que las interpreta con el sentido que tienen hoy.

\begin{exbox}
La tenencia de estupefacientes para consumo personal plantea la cuestión de qué
extensión darle a las expresiones del artículo 19 de la Constitución Nacional
según las cuales las acciones privadas no pueden ofender la moral ni las buenas
costumbres, ni perjudicar a terceros (desarrollado en profundidad en el
Capítulo~\ref{cap:p12}). Una lectura textualista acota la interpretación a la
literalidad de esas palabras; otras corrientes admiten un margen mayor para
proyectar la norma a los tiempos actuales.
\end{exbox}

\subsection{El originalismo}

\begin{defbox}[: Originalismo]
Corriente según la cual la interpretación debe responder a la fidelidad del texto,
pero además debe estar imbuida del período histórico en que la norma fue redactada:
las palabras tienen el sentido que tenían históricamente al momento de su sanción,
según la voluntad de quienes la dictaron.
\end{defbox}

El originalismo no es equivalente al textualismo, aunque ambas son corrientes
\textit{restrictivas} de la actividad judicial: buscan achicar el margen de
creación del juez, pero por vías distintas. El textualismo se apoya en el texto y
su significado actual; el originalismo agrega el peso histórico de la voluntad de
los constituyentes.

\begin{table}[H]
\centering
\caption{Textualismo y originalismo: dos variantes conservadoras}
\begin{tabularx}{\textwidth}{@{}lXX@{}}
\toprule
\textbf{Dimensión} & \textbf{Textualismo} & \textbf{Originalismo} \\
\midrule
Fuente primaria & Letra de la norma & Letra y contexto histórico de sanción \\
Significado de las palabras & Uso corriente y actual & Sentido al momento de la redacción \\
Rol del juez & Muy limitado por el texto & Limitado por texto e historia legislativa \\
Arraigo en Argentina & Moderado (doctrina de la Corte) & Débil \\
Arraigo en EE.\,UU. & Fuerte & Muy fuerte (jueces conservadores) \\
\bottomrule
\end{tabularx}
\end{table}

En el derecho argentino el originalismo tiene un arraigo más débil que en Estados
Unidos, donde la apelación a la voluntad de los ``padres fundadores'' es una
tradición consolidada. Un condicionante práctico explica en parte esa diferencia:
no existe una constancia documentada tan clara de los debates de la Asamblea
Constituyente de 1853 (la de 1860 está algo más documentada) como sí existe de la
Convención de Filadelfia de 1787.

\subsection{El pragmatismo o interpretación dinámica}

\begin{defbox}[: Pragmatismo]
La corriente más flexible de las tres. Prioriza mantener actualizados los
principios constitucionales esenciales por sobre la voluntad explícita del
legislador originario o la letra estricta, otorgando al intérprete un rol más
activo y creativo.
\end{defbox}

El pragmatismo es la corriente con mayor fertilidad interpretativa, porque no busca
su fundamento en el texto ni en la intención histórica, sino en mantener vigentes
los principios esenciales de la Constitución, reconociendo que los constituyentes
originales —aunque sabios— fueron ``producto de un tiempo que ya no existe''.

\subsection{Textualismo frente a no textualismo}

Vista en conjunto, la clasificación de la guía puede reconstruirse así:

\begin{table}[H]
\centering
\caption{Textualismo y no textualismo}
\begin{tabularx}{\textwidth}{@{}lX@{}}
\toprule
\textbf{Postura} & \textbf{Corrientes que la integran y su lógica común} \\
\midrule
Textualismo & El textualismo puro: la letra, en su sentido actual, es la única
fuente legítima y el límite del intérprete. \\
No textualismo & El originalismo (se aparta de la letra actual para atender a la
voluntad histórica) y el pragmatismo (se aparta tanto de la letra como de la
historia para atender a los fines y principios actuales). Ambas variantes admiten,
por caminos distintos, que el intérprete mire más allá del texto puro. \\
\bottomrule
\end{tabularx}
\end{table}

\begin{impbox}
No confundir ``no textualista'' con ``antitextual''. Ni el originalismo ni el
pragmatismo prescinden del texto: ambos parten de él, pero admiten —a diferencia
del textualismo puro— que otros elementos (la historia, los principios, la
finalidad) intervengan legítimamente en la interpretación final.
\end{impbox}

Como complemento transversal, la cátedra incorpora el \textbf{análisis económico
del derecho}: la idea de que todo gran diseño institucional tiene efectos
económicos concretos, y que cualquier análisis de un fallo o de una decisión
judicial debería tenerlos en cuenta, sin que ello convierta a la materia en una
disciplina económica.

\section{\textit{Holding} y \textit{obiter dictum} en el precedente judicial}

\begin{notebox}[Ampliación externa]
Los apuntes de la cátedra no traían un desarrollo propio de esta distinción. El
contenido que sigue se elaboró consultando doctrina procesal y constitucional
argentina especializada en la teoría del precedente (véase la bibliografía al
final del volumen).
\end{notebox}

Cuando un tribunal dicta sentencia, no todas las razones que expone en sus
considerandos tienen el mismo valor como precedente para casos futuros. La teoría
del precedente distingue entre lo que fue efectivamente decisivo para resolver el
caso y lo que el tribunal dijo ``de paso'', sin que fuera necesario para llegar a
esa decisión.

\begin{defbox}[: \textit{Holding} (o \textit{ratio decidendi})]
Es la regla de derecho que resulta estrictamente necesaria para resolver el caso
concreto: el razonamiento dirimente sobre cuya base el tribunal llega a la
decisión. Es la parte de la sentencia dotada de fuerza de precedente, la que
vincula a los tribunales inferiores y orienta la resolución de casos análogos
futuros.
\end{defbox}

\begin{defbox}[: \textit{Obiter dictum}]
Expresión latina que significa literalmente ``dicho de paso''. Son los argumentos,
reflexiones o afirmaciones que aparecen en la parte considerativa de una sentencia,
pero que no resultan imprescindibles para la decisión adoptada: acompañan,
corroboran o contextualizan el razonamiento central, sin ser su fundamento
dirimente. Por eso no sientan precedente obligatorio en sentido estricto, aunque
pueden tener fuerza persuasiva si provienen de un tribunal de gran prestigio o si,
por su claridad, son retomados en fallos posteriores.
\end{defbox}

\begin{table}[H]
\centering
\caption{\textit{Holding} y \textit{obiter dictum}}
\begin{tabularx}{\textwidth}{@{}lXX@{}}
\toprule
\textbf{Aspecto} & \textbf{\textit{Holding}} & \textbf{\textit{Obiter dictum}} \\
\midrule
Relación con el fallo & Es la razón que decide el caso & Es periférico: acompaña
sin decidir \\
Fuerza como precedente & Vinculante (en el sistema argentino, con matices: la
CSJN no está formalmente atada a sus propios precedentes, salvo sentencias
plenarias, pero su autoridad institucional genera un fuerte deber moral de
seguirlos) & Meramente persuasivo, según el prestigio del tribunal \\
Reversibilidad & Puede perder su carácter dirimente si otro argumento pasa a
sostener la decisión & Puede pasar a ser dirimente si, atacado con éxito el
argumento principal, el obiter alcanza por sí solo para sostener la decisión \\
\bottomrule
\end{tabularx}
\end{table}

\begin{exbox}
El ejemplo más célebre de esta distinción en el derecho comparado es la nota al
pie n.\textsuperscript{o} 4 del caso \textit{United States v.\ Carolene Products
Co.} (Corte Suprema de los Estados Unidos, 1938): un párrafo que, técnicamente, era
un \textit{obiter dictum} —no era necesario para resolver el caso, referido a la
regulación de la leche reconstituida—, pero que terminó siendo la base doctrinal
de todo el desarrollo posterior sobre el escrutinio judicial reforzado a favor de
las minorías discretas e insulares, como se retoma en el Capítulo~\ref{cap:p4}.
\end{exbox}

\begin{impbox}
Por qué importa distinguirlos al analizar un fallo. Cuando se estudia un
precedente para responder una pregunta de examen, no alcanza con citar ``lo que
dijo la Corte'': hay que identificar cuál fue el \textit{holding} —la regla que
efectivamente decidió el caso y que puede invocarse como precedente— y separar los
\textit{obiter dicta}, que sirven para comprender el contexto y el razonamiento del
tribunal, pero no pueden invocarse con el mismo peso normativo.
\end{impbox}

\section{Cuadro Cornell: interpretación y precedente}

\begin{longtable}{@{}p{0.38\textwidth}p{0.55\textwidth}@{}}
\toprule
\textbf{Preguntas / palabras clave} & \textbf{Notas / respuestas} \\
\midrule
\endhead
¿Cuál es el problema central de la interpretación constitucional? &
Que un texto fundacional, estable en el tiempo, debe aplicarse a una realidad
social que cambia mucho más rápido, sin que medie necesariamente una reforma. \\
¿Qué es el textualismo? & La corriente que da a la letra, en su sentido actual,
el rol de primera y casi excluyente fuente de interpretación, reduciendo al mínimo
la creatividad judicial. \\
¿En qué se diferencia el originalismo del textualismo? & El originalismo agrega el
peso de la voluntad histórica de los constituyentes; el textualismo se limita al
significado actual de las palabras. \\
¿Qué caracteriza al pragmatismo? & Prioriza mantener vigentes los principios
esenciales de la Constitución por sobre la letra o la historia, con mayor
fertilidad para la reinterpretación. \\
¿Cómo se agrupan estas corrientes en textualismo y no textualismo? & El
textualismo por un lado; el originalismo y el pragmatismo, por vías distintas,
como variantes del no textualismo, en tanto ambos admiten mirar más allá de la
letra pura. \\
¿Qué es el \textit{holding}? & La razón estrictamente necesaria para decidir el
caso: la regla que sienta precedente. \\
¿Qué es el \textit{obiter dictum}? & Lo dicho ``de paso'' en la sentencia, que no
fue necesario para la decisión y que solo tiene fuerza persuasiva. \\
¿Por qué importa la nota al pie n.\textsuperscript{o} 4 de \textit{Carolene
Products}? & Porque, siendo técnicamente un \textit{obiter dictum}, se convirtió en
el fundamento doctrinal del escrutinio judicial reforzado a favor de minorías
discretas e insulares. \\
\midrule
\multicolumn{2}{p{0.95\textwidth}}{\textbf{Resumen:} frente a un texto que
envejece, tres corrientes disputan su interpretación —textualismo, originalismo y
pragmatismo—, agrupables en la dicotomía textualismo/no textualismo según su apego
a la letra. Al analizar cualquier precedente judicial, es indispensable distinguir
el \textit{holding} —la razón dirimente, dotada de fuerza de precedente— del
\textit{obiter dictum} —lo dicho de paso, con valor meramente persuasivo—, aunque
un \textit{obiter} especialmente influyente, como la nota 4 de \textit{Carolene
Products}, puede terminar orientando toda una línea jurisprudencial posterior.} \\
\bottomrule
\end{longtable}

% ============================================================
% PARTE II - CONTROL DE CONSTITUCIONALIDAD
% ============================================================
\part{Control de constitucionalidad}

\chapter{Concepto y alcance del control de constitucionalidad. El papel de los tribunales y de la Corte Suprema. La supremacía de la Constitución Nacional}
\label{cap:p2}

\section{El punto de partida: la supremacía constitucional}

\begin{defbox}[: Supremacía constitucional]
Principio, consagrado en el artículo 31 de la Constitución Nacional, según el cual
la Constitución Nacional, las leyes que en su consecuencia dicte el Congreso y los
tratados con las potencias extranjeras son la ley suprema de la Nación, y las
autoridades de cada provincia están obligadas a conformarse a ella, no obstante
cualquier disposición en contrario que contengan las leyes o constituciones
provinciales.
\end{defbox}

La supremacía constitucional implica una jerarquía: ninguna norma inferior —ley,
decreto, ordenanza, sentencia— puede contradecir a la Constitución. Pero afirmar
esa jerarquía en abstracto no basta: hace falta un mecanismo que, en cada caso
concreto, permita verificar si una norma inferior efectivamente la respeta, y que
prive de efectos a la que no lo haga. Ese mecanismo es el \textbf{control de
constitucionalidad}.

\begin{defbox}[: Control de constitucionalidad]
Mecanismo institucional por el cual un órgano del Estado —en la Argentina,
cualquier juez, con la Corte Suprema como intérprete final— verifica si una norma o
un acto de los poderes públicos se ajusta a la Constitución y, en caso de
contradicción, declara su invalidez para el caso concreto, privándola de
aplicación.
\end{defbox}

\section{El alcance del control: qué puede controlarse y quién lo hace}

En el sistema argentino, cualquier juez llamado a resolver un conflicto puede,
siempre que se trate de una cuestión justiciable —es decir, de un caso o
controversia concreto, no de una consulta abstracta—, ejercer el control de
constitucionalidad y, como \textit{última ratio} del ordenamiento, invalidar una
norma emanada del Congreso o del Poder Ejecutivo: una ley, un decreto o un decreto
de necesidad y urgencia. Este es el llamado \textbf{sistema difuso}, que se
desarrolla en detalle en el Capítulo~\ref{cap:p3}.

El control no se limita a las leyes: alcanza también a los decretos, a los actos
administrativos y, en general, a cualquier acto de los poderes públicos que pueda
colisionar con la Constitución. Lo que lo activa siempre es un caso judicial
concreto: sin un conflicto de intereses entre partes que deba resolverse, no hay
control de constitucionalidad posible, porque el Poder Judicial argentino no
resuelve consultas abstractas sobre la validez de las normas.

\section{El papel de los tribunales inferiores y el rol final de la Corte Suprema}

Si todos los jueces pueden declarar la inconstitucionalidad de una norma, es
imprescindible que exista un órgano que ordene y unifique esas decisiones, para
evitar el caos de sentencias contradictorias. Ese rol lo cumple la \textbf{Corte
Suprema de Justicia de la Nación}, como intérprete final y máximo de la
Constitución. La Corte no es el único tribunal habilitado para ejercer el control
—cualquier juez de primera instancia puede hacerlo—, pero sí es la última instancia
a la que, mediante el recurso extraordinario federal (desarrollado en el
Capítulo~\ref{cap:p8}), puede llegar cualquier cuestión constitucional, para que
fije un criterio uniforme y definitivo.

\begin{impbox}
El sistema difuso necesita un principio ordenador. Sin un órgano que unifique el
criterio, el sistema difuso caería en decisiones caóticas y contradictorias. Ese
rol ordenador lo cumple la Corte Suprema, no porque monopolice el control, sino
porque es el intérprete final ante quien todas las líneas de interpretación
terminan convergiendo.
\end{impbox}

\section{Los fallos fundacionales: Marbury v.\ Madison y Sojo}

El origen histórico de esta atribución judicial —que ninguna Constitución escrita
establece de manera expresa y detallada— se remonta a dos fallos fundacionales,
uno norteamericano y uno argentino, que responden a la misma pregunta: quién tiene
la última palabra sobre el significado de la Constitución.

\subsection{Marbury v.\ Madison (Corte Suprema de los Estados Unidos, 1803)}

El trasfondo político del caso fue el siguiente: el presidente saliente John Adams,
tras perder la elección frente a Thomas Jefferson, había nombrado a 42 jueces que
todavía no habían asumido sus cargos. Jefferson, ya electo, se negó a completar esos
nombramientos y sostuvo que él mismo podía interpretar la Constitución sin
someterse al criterio de los tribunales. William Marbury, uno de los jueces
nombrados, reclamó directamente ante la Corte Suprema que se lo pusiera en
funciones.

El presidente de la Corte, John Marshall, resolvió el caso con un razonamiento de
tres preguntas:

\begin{table}[H]
\centering
\caption{La ``lógica de Marshall'' en \textit{Marbury v.\ Madison}}
\begin{tabularx}{\textwidth}{@{}p{0.05\textwidth}Xl@{}}
\toprule
\textbf{N.\textsuperscript{o}} & \textbf{Pregunta} & \textbf{Respuesta} \\
\midrule
1.\textsuperscript{a} & ¿Existió el acto de nombramiento? & Sí \\
2.\textsuperscript{a} & ¿Tiene Marbury un derecho y acudió al lugar adecuado? & Sí \\
3.\textsuperscript{a} & ¿Puede el Poder Judicial ordenar, mediante un
\textit{mandamus}, que se lo ponga en el cargo? & No: invadiría al Ejecutivo \\
\bottomrule
\end{tabularx}
\end{table}

Ante la tercera pregunta la Corte se autorrestringió —declaró inconstitucional la
sección de la \textit{Judiciary Act} que le daba competencia originaria para
emitir el mandamus— y, paradójicamente, ganó más: quedó consagrada como el único
órgano capacitado para interpretar y decir lo que es la ley, incluso frente al
Presidente. De allí la frase célebre del fallo: ``el gobierno de los Estados
Unidos es un gobierno de leyes y no de hombres''.

\begin{impbox}
``No tenemos ni la bolsa ni la espada''. El Poder Judicial no es el más fuerte
de los tres poderes, sino el más débil: no controla los fondos públicos (la
``bolsa'', en manos del Legislativo) ni el monopolio de la fuerza (la ``espada'',
en manos del Ejecutivo). Su único recurso es la autoridad de decir lo que la
Constitución significa. Esta debilidad institucional es, precisamente, la que
alimenta la objeción contramayoritaria que se estudia en el
Capítulo~\ref{cap:p4}: un poder sin respaldo electoral directo termina, sin
embargo, invalidando las decisiones de los poderes que sí lo tienen.
\end{impbox}

\subsection{Sojo c/ Poder Ejecutivo Nacional (CSJN, 1887)}

En la Argentina, apenas veinticinco años después de sancionada la Constitución de
1853/60, se presentó un caso análogo. Eduardo Sojo, un dibujante que publicó una
caricatura política satírica, fue detenido sin juicio previo por orden del
Congreso. Su abogado planteó un habeas corpus directamente ante la Corte Suprema,
en lugar de ante un juez de primera instancia.

La Corte interpretó los artículos 116 y 117 de la Constitución Nacional
(entonces 100 y 101): el artículo 116 establece la regla general, la competencia
por apelación, agotando las instancias desde la más baja; el artículo 117
establece la excepción, la competencia originaria y exclusiva de la Corte,
reservada —bajo un \textit{numerus clausus} taxativo— a los casos que involucran
embajadores, ministros y cónsules extranjeros, y a las causas en que una provincia
sea parte. Como Sojo no encuadraba en ninguno de esos supuestos, la Corte remitió
la causa al juez de primera instancia, quien hizo lugar al habeas corpus y liberó
a Sojo.

\begin{table}[H]
\centering
\caption{Marbury y Sojo: comparación}
\begin{tabularx}{\textwidth}{@{}lXX@{}}
\toprule
& \textbf{Marbury v.\ Madison} & \textbf{Sojo c/ Poder Ejecutivo} \\
\midrule
Cuestión de fondo & Control de constitucionalidad & Jurisdicción y competencia, no
inconstitucionalidad \\
Vía elegida & Directamente ante la Corte & Directamente ante la Corte \\
Resultado & La Corte se consolida como intérprete final de la Constitución & La
Corte remite el caso a primera instancia; Sojo es liberado allí \\
\bottomrule
\end{tabularx}
\end{table}

Aunque los resultados procesales fueron distintos, el efecto sistémico es el
mismo: en ambos casos queda consolidada la idea de que son los tribunales —y en
última instancia la Corte Suprema— quienes dicen qué significa la Constitución, y
que ese poder de control corresponde, en el sistema argentino, a todos los jueces
y no solo al tribunal cimero (control difuso, Capítulo~\ref{cap:p3}).

\section{Cuadro Cornell: concepto, alcance y supremacía}

\begin{longtable}{@{}p{0.38\textwidth}p{0.55\textwidth}@{}}
\toprule
\textbf{Preguntas / palabras clave} & \textbf{Notas / respuestas} \\
\midrule
\endhead
¿Qué es la supremacía constitucional? & El principio del artículo 31 CN según el
cual la Constitución, las leyes que en su consecuencia se dicten y los tratados
son la ley suprema de la Nación, por sobre cualquier norma provincial o inferior. \\
¿Qué es el control de constitucionalidad? & El mecanismo por el cual un juez, en
un caso concreto, verifica si una norma respeta la Constitución y, si no lo hace,
la declara inválida para ese caso. \\
¿Qué hace falta para que proceda el control? & Un caso o controversia judicial
concreto y justiciable; no procede como consulta abstracta. \\
¿Quién ejerce el control en Argentina? & Cualquier juez, en cualquier proceso
(sistema difuso), con la Corte Suprema como intérprete final y ordenador del
sistema. \\
¿Qué estableció \textit{Marbury v.\ Madison}? & Que el Poder Judicial es el único
órgano capacitado para decir qué significa la Constitución, incluso frente al
Presidente, aunque el Poder Judicial ``no tiene ni la bolsa ni la espada''. \\
¿Qué resolvió \textit{Sojo}? & Que el caso no era de inconstitucionalidad sino de
competencia: la vía originaria ante la Corte es excepcional y taxativa (art.\ 117
CN); la regla es la apelación (art.\ 116 CN). \\
\midrule
\multicolumn{2}{p{0.95\textwidth}}{\textbf{Resumen:} la supremacía constitucional
(art.\ 31 CN) exige un mecanismo que la haga efectiva en cada caso: el control de
constitucionalidad, ejercido en la Argentina por todos los jueces (sistema
difuso) y ordenado, en última instancia, por la Corte Suprema. \textit{Marbury} y
\textit{Sojo} son los fallos fundacionales que consolidaron esta atribución
judicial en Estados Unidos y en la Argentina, respectivamente.} \\
\bottomrule
\end{longtable}


\chapter{Control difuso y concentrado de constitucionalidad. Diferencias, ventajas y desventajas; contexto jurídico institucional de cada sistema}
\label{cap:p3}

\section{El sistema difuso}

\begin{defbox}[: Control de constitucionalidad difuso]
Sistema, adoptado por la Argentina y por los Estados Unidos, en el que todos los
jueces —y no solo la Corte Suprema— tienen el poder de declarar inconstitucional
una norma en un caso concreto, en la medida en que sea conducente para resolver el
conflicto. No existe un tribunal con competencia exclusiva y excluyente para
resolver esa cuestión, como sí sucede en otros ordenamientos.
\end{defbox}

En el sistema argentino, cualquier juez, al resolver un caso, puede invalidar
—como última ratio— una ley, un decreto o un decreto de necesidad y urgencia que
contradiga la Constitución. Ese poder no está concentrado en un tribunal especial:
está \textit{difundido} entre toda la magistratura.

\section{El sistema concentrado}

En el sistema continental europeo —por ejemplo, en Italia, Alemania y España—
existen \textbf{tribunales constitucionales}, que ejercen en forma exclusiva y
excluyente el control de constitucionalidad. Un juez ordinario, en esos sistemas,
no puede por sí mismo declarar inválida una ley: debe plantear la cuestión ante el
tribunal especializado, que es el único habilitado para resolverla, generalmente
con efectos generales (\textit{erga omnes}) hacia el futuro.

\section{Comparación: diferencias, ventajas y desventajas}

\begin{table}[H]
\centering
\caption{Sistema difuso y sistema concentrado}
\begin{tabularx}{\textwidth}{@{}lXX@{}}
\toprule
\textbf{Característica} & \textbf{Sistema difuso (Argentina / EE.\,UU.)} &
\textbf{Sistema concentrado (España / Alemania / Italia)} \\
\midrule
¿Quién ejerce el control? & Cualquier juez, en cualquier proceso & Un tribunal
constitucional especializado y exclusivo \\
Acceso del justiciable & Directo, en el mismo proceso & Indirecto: se remite al
tribunal especializado \\
Alcance de la decisión & Para el caso concreto (\textit{inter partes}) & Efecto
general (\textit{erga omnes}) \\
Riesgo principal & Sentencias contradictorias entre juzgados & Menor acceso
directo para justiciables sin recursos para litigar hasta esa instancia \\
Carácter & Más democrático y accesible & Más especializado, más previsible en su
efecto general \\
\bottomrule
\end{tabularx}
\end{table}

\subsection{Ventajas del sistema difuso}

El sistema difuso, con sus imperfecciones, resulta más democrático porque está al
alcance de todos: cualquier abogado puede plantear la invalidez de una norma en
cualquier proceso; el justiciable no debe esperar la intervención de un tribunal
lejano y especializado, porque la cuestión se resuelve en el mismo expediente; y
el paso por las sucesivas instancias recursivas enriquece el debate y permite una
maduración progresiva de los argumentos antes de llegar, eventualmente, a la Corte
Suprema.

\subsection{Costos y desventajas del sistema difuso}

Como contrapartida, el sistema difuso genera un plazo de incertidumbre: pueden
existir, durante un tiempo, sentencias contrapuestas sobre la misma norma; una ley
general puede quedar puesta en duda o inaplicada en algunos casos y no en otros,
generando la situación —a veces contraintuitiva pero real— de que la misma norma
se aplique para algunos justiciables y no para otros; y pueden generarse tensiones
entre los poderes del Estado mientras la cuestión no queda zanjada definitivamente.
Estos son los llamados \textbf{costos de transacción} del sistema difuso.

\begin{impbox}
El sistema difuso necesita un principio ordenador. Los costos de transacción
recién descriptos explican por qué el sistema difuso no podría funcionar sin la
Corte Suprema actuando como intérprete final y unificador de criterio, como se
explicó en el Capítulo~\ref{cap:p2}: sin ese rol ordenador, el sistema caería en
decisiones caóticas y contradictorias.
\end{impbox}

\subsection{Contexto jurídico-institucional de cada sistema}

La diferencia entre ambos sistemas no es casual, sino que responde a tradiciones
jurídicas distintas. El sistema difuso nace en el \textit{common law}
norteamericano, donde los jueces gozan históricamente de un fuerte poder creador y
donde la fuerza del precedente judicial (\textit{stare decisis}) permite, a pesar
de la dispersión del control, mantener cierta uniformidad interpretativa a través
del tiempo. La Argentina, país de tradición continental en su derecho de fondo pero
que copió el modelo de control judicial norteamericano, combina ambos elementos: un
Código Civil y Comercial y un Código Penal únicos para todo el país (a diferencia
de Estados Unidos, donde cada estado tiene su propia legislación de fondo), con un
control de constitucionalidad difuso a la manera estadounidense.

El sistema concentrado, en cambio, nace en la tradición continental europea de
posguerra (el modelo austríaco de Hans Kelsen, adoptado luego por Alemania,
Italia y España), en sistemas donde el Poder Judicial ordinario, por tradición,
tenía un rol más restringido frente a la ley del Parlamento, y donde se prefirió
crear un órgano especial, con una legitimidad y una composición particulares
(a menudo con participación de otros poderes en la designación de sus miembros),
para concentrar la tarea, sensible desde el punto de vista institucional, de
invalidar leyes con efectos generales.

\section{Cuadro Cornell: control difuso y concentrado}

\begin{longtable}{@{}p{0.38\textwidth}p{0.55\textwidth}@{}}
\toprule
\textbf{Preguntas / palabras clave} & \textbf{Notas / respuestas} \\
\midrule
\endhead
¿Qué es el sistema difuso? & Aquel en que cualquier juez puede declarar la
inconstitucionalidad de una norma en el caso concreto que está resolviendo. \\
¿Qué es el sistema concentrado? & Aquel en que solo un tribunal constitucional
especializado y exclusivo puede declarar la inconstitucionalidad, generalmente con
efectos \textit{erga omnes}. \\
¿Cuál es la principal ventaja del sistema difuso? & Su accesibilidad democrática:
cualquier justiciable puede plantear la cuestión en el mismo proceso, sin esperar
a un tribunal especial. \\
¿Cuál es su principal costo? & La posibilidad de sentencias contradictorias y de
un período de incertidumbre mientras la cuestión no llega a la Corte Suprema. \\
¿De qué tradición proviene cada sistema? & El difuso, del \textit{common law}
norteamericano; el concentrado, de la tradición continental europea de posguerra
(modelo kelseniano). \\
\midrule
\multicolumn{2}{p{0.95\textwidth}}{\textbf{Resumen:} el sistema difuso argentino
—cualquier juez, efecto \textit{inter partes}— es más accesible y democrático,
pero exige un principio ordenador (la Corte Suprema) para evitar sentencias
contradictorias; el sistema concentrado europeo —un tribunal constitucional
exclusivo, efecto \textit{erga omnes}— gana en uniformidad y previsibilidad, pero
resulta menos accesible en forma directa para el justiciable común.} \\
\bottomrule
\end{longtable}

\chapter{El Poder Judicial como poder contramayoritario}
\label{cap:p4}

\begin{notebox}[Ampliación externa]
Los apuntes de la cátedra no traían un desarrollo propio de este punto, aunque el
cronograma de la Comisión 603 lo incluye expresamente (clase del 31/08, con
bibliografía de Ely y el fallo \textit{United States v.\ Carolene Products Co.}).
El desarrollo que sigue se elaboró consultando doctrina constitucional especializada
en la llamada ``dificultad contramayoritaria'' (véase la bibliografía final).
\end{notebox}

\section{El problema: jueces no electos que invalidan leyes de representantes electos}

Como se vio en el Capítulo~\ref{cap:p2}, a partir de \textit{Marbury v.\ Madison} y
de \textit{Sojo} quedó consolidada la atribución de los jueces de declarar
inválida una ley del Congreso o un acto del Poder Ejecutivo quc contradiga la
Constitución. Pero esa atribución plantea una tensión que la teoría democrática no
puede ignorar: los jueces no son elegidos por el voto popular —o lo son solo de
manera muy mediata, a través de mecanismos de designación en los que interviene el
Poder Ejecutivo y el Senado—, y sin embargo tienen la última palabra para dejar sin
efecto una ley sancionada por los representantes que sí fueron elegidos
directamente por el pueblo.

\begin{defbox}[: Dificultad contramayoritaria (\textit{countermajoritarian
difficulty})]
Expresión acuñada por el jurista norteamericano Alexander Bickel en su obra
\textit{The Least Dangerous Branch} (1962) para describir la tensión estructural
entre el control judicial de constitucionalidad y el principio democrático: cuando
la Corte Suprema declara inconstitucional un acto legislativo, frustra la voluntad
de los representantes elegidos por la mayoría, ejerciendo así un poder que, en ese
punto concreto, opera \textit{contra} la mayoría y sin el respaldo directo del voto
popular.
\end{defbox}

Bickel no concluye que el control judicial deba abandonarse: sostiene que se trata
de una \textit{dificultad} que debe administrarse con prudencia, y desarrolla la
idea de las ``virtudes pasivas'': mecanismos procesales (por ejemplo, no tratar
cuestiones políticas no justiciables, o esperar a que un caso esté suficientemente
``maduro'') mediante los cuales la propia Corte puede evitar pronunciarse quc
sobre cuestiones que la enfrentarían innecesariamente con las mayorías políticas,
reservando su intervención más fuerte para los casos en que resulte
verdaderamente indispensable.

\section{La respuesta de Ely: la teoría del refuerzo de representación}

\begin{notebox}[Ampliación externa]
El cronograma indica como bibliografía especial de este punto a John Hart Ely,
\textit{Democracia y desconfianza}. El desarrollo que sigue sintetiza esa obra a
partir de fuentes académicas secundarias, dado que los apuntes de clase no
registraban su contenido.
\end{notebox}

El constitucionalista norteamericano John Hart Ely propuso, en \textit{Democracy
and Distrust} (1980), una respuesta a la dificultad contramayoritaria que no niega
el problema, pero busca reformularlo: si el control judicial se ejerciera para
imponer los valores sustantivos preferidos por los jueces (qué es una vida buena,
qué política económica es correcta), efectivamente usurparía el lugar de las
mayorías democráticas. Pero si el control judicial se concentra, en cambio, en
\textbf{garantizar que el propio proceso democrático funcione correctamente} —que
los canales de participación y de cambio político estén abiertos, y que las
minorías no queden sistemáticamente excluidas de ellos—, entonces los jueces no
están sustituyendo a la democracia: están \textit{reforzándola}.

\begin{defbox}[: Teoría del refuerzo de representación]
Propuesta de John Hart Ely según la cual el control judicial de constitucionalidad
se legitima, frente a la objeción contramayoritaria, cuando se dirige a
desbloquear o a corregir fallas del propio proceso político-representativo: (i)
cuando quienes están en el poder distorsionan los canales de cambio político para
perpetuarse (por ejemplo, restringiendo el voto o la competencia electoral), o (ii)
cuando una mayoría sistemáticamente perjudica a una minoría discreta e insular
sin que exista un mecanismo político realista para que esa minoría se defienda a
sí misma. En ambos casos, el juez no reemplaza el juicio de valor de la mayoría:
repara una falla en el mecanismo por el cual esa mayoría se forma y decide.
\end{defbox}

\subsection{La nota al pie n.\textsuperscript{o} 4 de \textit{Carolene Products}}

El punto de partida doctrinal de Ely es un párrafo del fallo \textit{United
States v.\ Carolene Products Co.} (Corte Suprema de EE.\,UU., 1938), redactado por
el juez Harlan Fiske Stone: técnicamente un \textit{obiter dictum}
(véase Capítulo~\ref{cap:p1}), pero que se convirtió en una de las notas al pie
más influyentes del derecho constitucional comparado. Ese párrafo sugiere que la
presunción habitual de constitucionalidad de las leyes debe ceder a un escrutinio
judicial más estricto en tres supuestos: cuando la norma restringe directamente el
propio proceso político (por ejemplo, el derecho al voto); cuando distorsiona los
canales de reforma legislativa que permitirían corregir una legislación
indeseable por vías democráticas; o cuando perjudica a lo que la nota llama
\textbf{``minorías discretas e insulares''}.

\begin{defbox}[: Minorías discretas e insulares]
Grupos minoritarios —definidos, por ejemplo, por su religión, su origen étnico o
nacional, o rasgos análogos— que se encuentran aislados del proceso político
ordinario: por su tamaño reducido, por los prejuicios que enfrentan o por la
dificultad para construir alianzas con otros grupos, no logran hacerse oír ni
defender sus intereses a través de los mecanismos habituales de la democracia
mayoritaria (el voto, el lobby, la formación de coaliciones). Frente a ellas, el
proceso político en el que habitualmente se confía para proteger a las minorías
deja de funcionar, y se justifica una revisión judicial más exigente.
\end{defbox}

\begin{impbox}
Relación con las categorías sospechosas. Esta misma lógica reaparece, en el
derecho argentino, bajo la doctrina de las \textbf{categorías sospechosas}
(desarrollada en el Capítulo~\ref{cap:p13} a propósito de la libertad religiosa):
cuando una norma distingue por motivos como la religión, el origen nacional o
rasgos análogos, se invierte la presunción de legitimidad de la ley, exigiéndose
al Estado que demuestre su razonabilidad. Es la misma preocupación por las
minorías discretas e insulares, trasladada a una técnica de control judicial
concreta.
\end{impbox}

\section{El poder judicial y los procesos políticos: \textit{Baker v.\ Carr}}

El cronograma vincula esta unidad con el fallo \textit{Baker v.\ Carr} (Corte
Suprema de los Estados Unidos, 1962), en el que la Corte admitió, por primera vez,
que los tribunales federales podían revisar judicialmente la manera en que los
estados trazaban sus distritos electorales (la llamada ``malapportionment''),
rechazando la idea de que se tratara de una ``cuestión política no justiciable''
ajena al control judicial. El fallo ilustra exactamente el segundo tipo de falla
que identifica Ely: cuando quienes están en el poder distorsionan las propias
reglas del proceso electoral para perpetuarse, ese no es un tema que la mayoría
pueda corregirse a sí misma por la vía ordinaria —porque son, precisamente, los
beneficiarios de la distorsión quienes controlan el proceso legislativo—, y allí
se justifica la intervención judicial: no para imponer un resultado electoral
determinado, sino para reabrir el propio canal de cambio político.

\section{Síntesis: cómo se articula la objeción contramayoritaria con su respuesta}

\begin{table}[H]
\centering
\caption{La dificultad contramayoritaria y la respuesta de Ely}
\begin{tabularx}{\textwidth}{@{}lX@{}}
\toprule
\textbf{Elemento} & \textbf{Contenido} \\
\midrule
Objeción (Bickel) & El control judicial de constitucionalidad es una fuerza
contramayoritaria: jueces no electos invalidan leyes de representantes electos. \\
Respuesta (Ely) & El control judicial se legitima cuando repara fallas del propio
proceso democrático: canales de participación bloqueados, o perjuicio
sistemático a minorías discretas e insulares que no pueden defenderse por la vía
política ordinaria. \\
Fundamento doctrinal & Nota al pie n.\textsuperscript{o} 4 de \textit{Carolene
Products} (1938), técnicamente un \textit{obiter dictum} que se volvió
doctrina central. \\
Aplicación jurisprudencial & \textit{Baker v.\ Carr} (1962): revisión judicial de
la distribución de distritos electorales, por afectar el propio proceso de
representación. \\
Proyección en el derecho argentino & La doctrina de las categorías sospechosas
(Capítulo~\ref{cap:p13}), que invierte la carga de la prueba de razonabilidad
cuando una norma distingue por religión, origen u otros rasgos sensibles. \\
\bottomrule
\end{tabularx}
\end{table}

\begin{exambox}
El Poder Judicial es contramayoritario porque sus integrantes no son elegidos
por el voto popular —o lo son solo de manera muy indirecta— y, sin embargo, pueden
invalidar leyes sancionadas por los representantes de la mayoría (Bickel). Esta
tensión no elimina la legitimidad del control judicial, pero exige justificarlo:
la respuesta más influyente es la de Ely, para quien los jueces refuerzan —en
lugar de reemplazar— a la democracia cuando intervienen para destrabar el propio
proceso político (por ejemplo, frente a distorsiones electorales, como en
\textit{Baker v.\ Carr}) o para proteger a minorías discretas e insulares que no
logran hacerse oír por los canales políticos ordinarios (nota 4 de
\textit{Carolene Products}). En ambos supuestos, el juez no impone su propia
preferencia de valores: corrige una falla del mecanismo por el cual la propia
mayoría se forma y decide.
\end{exambox}

\section{Cuadro Cornell: el Poder Judicial contramayoritario}

\begin{longtable}{@{}p{0.38\textwidth}p{0.55\textwidth}@{}}
\toprule
\textbf{Preguntas / palabras clave} & \textbf{Notas / respuestas} \\
\midrule
\endhead
¿Qué es la dificultad contramayoritaria? & La tensión, señalada por Bickel, entre
un Poder Judicial no electo y su facultad de invalidar leyes dictadas por
representantes elegidos por la mayoría. \\
¿Cómo responde Ely a esta dificultad? & Sosteniendo que el control judicial se
legitima cuando refuerza el propio proceso democrático: destrabando canales de
participación bloqueados o protegiendo a minorías discretas e insulares. \\
¿Qué es la nota 4 de \textit{Carolene Products}? & Un \textit{obiter dictum} de
1938 que enumeró los supuestos en que se justifica un escrutinio judicial más
estricto, y que se convirtió en el fundamento doctrinal de la teoría de Ely. \\
¿Qué son las minorías discretas e insulares? & Grupos minoritarios aislados del
proceso político ordinario, que no logran defender sus intereses por los
mecanismos habituales de la democracia mayoritaria. \\
¿Qué mostró \textit{Baker v.\ Carr}? & Que los tribunales pueden revisar la
distribución de distritos electorales, porque esa distorsión no puede corregirse
por la vía política ordinaria cuando son sus propios beneficiarios quienes
controlan el proceso legislativo. \\
\midrule
\multicolumn{2}{p{0.95\textwidth}}{\textbf{Resumen:} el Poder Judicial es
contramayoritario porque invalida, sin mandato electoral directo, decisiones de
los poderes representativos. Bickel plantea el problema; Ely ofrece la respuesta
más influyente: el control judicial se legitima cuando repara fallas del propio
proceso democrático —participación bloqueada o minorías discretas e insulares
desprotegidas—, en lugar de imponer valores sustantivos ajenos a la decisión
mayoritaria.} \\
\bottomrule
\end{longtable}

% ============================================================
% PARTE III - EL RECURSO EXTRAORDINARIO FEDERAL
% ============================================================
\part{El recurso extraordinario federal}

\chapter{Cuestión federal simple o compleja (directa e indirecta)}
\label{cap:p5}

\section{Ubicación del tema: el corazón del recurso extraordinario}

El recurso extraordinario federal (REF) es la vía por la cual una cuestión
constitucional llega, en última instancia, a la Corte Suprema. Su régimen general
—requisitos comunes y propios, sentencia definitiva, plazos— se desarrolla en el
Capítulo~\ref{cap:p8}. Pero antes de entrar en esos requisitos, es imprescindible
comprender su núcleo conceptual: la existencia de una \textbf{cuestión federal}.

\begin{impbox}
El corazón de los requisitos del recurso extraordinario está en la existencia de
una cuestión federal. El recurso procede formalmente contra sentencias definitivas
o asimilables (Capítulo~\ref{cap:p8}), pero, desde el punto de vista sustancial,
procede siempre que exista una controversia jurídica en la que esté en juego el
alcance de una norma federal o un conflicto entre normas de distinta jerarquía.
\end{impbox}

La Ley 48 establece la distinción, hoy clásica, entre cuestión federal
\textbf{simple} y \textbf{compleja}, y esta última se subdivide, a su vez, en
\textbf{directa} e \textbf{indirecta}.

\section{Cuestión federal simple}

\begin{defbox}[: Cuestión federal simple]
Controversia en la que está en juego la interpretación de la ley, de la
Constitución o de un tratado: se discute cuál es el verdadero alcance de una norma
federal, sin que exista un conflicto entre dos normas de distinta jerarquía.
\end{defbox}

\begin{exbox}
Si para resolver un caso es necesario desentrañar qué alcance corresponde darle al
concepto de ``indemnización plena'', previsto en una norma federal, se trata de una
cuestión federal simple: no hay dos normas en pugna, sino una sola cuyo sentido
exacto debe determinarse.
\end{exbox}

\section{Cuestión federal compleja: directa e indirecta}

\begin{defbox}[: Cuestión federal compleja]
A diferencia de la simple, en la cuestión federal compleja existe un
\textbf{conflicto normativo} entre dos o más normas.
\end{defbox}

\begin{table}[H]
\centering
\caption{Cuestión federal compleja directa e indirecta}
\begin{tabularx}{\textwidth}{@{}lXX@{}}
\toprule
\textbf{Tipo} & \textbf{Normas en conflicto} & \textbf{Ejemplo} \\
\midrule
Compleja directa & La Constitución Nacional y una norma de rango inferior & Una
ley o un decreto que contradicen directamente una disposición constitucional \\
Compleja indirecta & Dos o más normas inferiores a la Constitución, entre sí &
Una norma provincial que contradice una ley del Congreso, o un tratado sin
jerarquía constitucional, sin que la Constitución esté directamente en pugna \\
\bottomrule
\end{tabularx}
\end{table}

Un punto que suele generar confusión es el de los tratados internacionales. Los
tratados con jerarquía constitucional (art.\ 75, inc.\ 22 CN) forman parte de la
Constitución a estos efectos: si la controversia involucra a uno de esos tratados
frente a una norma inferior, se trata de una cuestión federal \textbf{compleja
directa}. Si, en cambio, involucra a un tratado sin jerarquía constitucional
—por ejemplo, un tratado del Mercosur— frente a una ley nacional, se trata de una
cuestión federal \textbf{compleja indirecta}, porque ambas normas son inferiores a
la Constitución y ninguna de ellas es, en sí misma, la Constitución.

\section{Elementos adicionales: decisión contraria y relación directa e inmediata}

Además de encuadrar el tipo de cuestión federal, dos requisitos adicionales
resultan indispensables:

\begin{defbox}[: Decisión contraria al derecho federal invocado]
Para que el recurso extraordinario sea procedente, la decisión impugnada debe ser
contraria al derecho federal invocado por quien recurre. Si la sentencia previa
resolvió aplicando correctamente la prelación constitucional (por ejemplo, haciendo
prevalecer una norma federal sobre una provincial, tal como la propia Constitución
lo exige), no hay cuestión constitucional en juego, porque la sentencia no es
contraria al derecho federal: lo aplicó correctamente.
\end{defbox}

\begin{defbox}[: Relación directa e inmediata]
La resolución de la cuestión federal debe tener una relación directa con la
decisión del pleito; no puede tratarse de una controversia meramente conjetural o
académica, desconectada del resultado concreto del proceso.
\end{defbox}

\begin{impbox}
¿Es obligatorio clasificar la cuestión federal? No es un requisito intrínseco de
admisibilidad etiquetar la cuestión como simple, compleja, directa o indirecta:
esa clasificación es una herramienta argumentativa —útil para ordenar el
planteo— pero lo verdaderamente indispensable es demostrar con precisión cuál es
el derecho federal vulnerado. Las categorías, además, no son puras ni excluyentes:
un mismo caso puede combinar, en parte, una controversia sobre el alcance de una
norma y, en parte, una colisión normativa propiamente dicha.
\end{impbox}

\section{Cuadro Cornell: cuestión federal}

\begin{longtable}{@{}p{0.38\textwidth}p{0.55\textwidth}@{}}
\toprule
\textbf{Preguntas / palabras clave} & \textbf{Notas / respuestas} \\
\midrule
\endhead
¿Qué es una cuestión federal simple? & La controversia sobre el alcance de una
norma federal, sin colisión con otra norma. \\
¿Qué diferencia a la cuestión federal compleja directa de la indirecta? & En la
directa, la Constitución colisiona con una norma inferior; en la indirecta,
colisionan dos o más normas inferiores entre sí. \\
¿Qué rol cumplen los tratados con jerarquía constitucional en esta clasificación?
& Se equiparan a la Constitución: si están en pugna con una norma inferior, la
cuestión es compleja directa. \\
¿Qué es la ``decisión contraria al derecho federal invocado''? & El requisito de
que la sentencia impugnada haya resuelto en contra del derecho federal que se
invoca, y no simplemente aplicado correctamente la prelación normativa. \\
¿Es obligatorio clasificar la cuestión federal? & No; lo esencial es demostrar
cuál es el derecho federal vulnerado. \\
\midrule
\multicolumn{2}{p{0.95\textwidth}}{\textbf{Resumen:} el núcleo del recurso
extraordinario es la cuestión federal: simple cuando se discute la
interpretación de una norma federal; compleja —directa o indirecta— cuando hay
colisión normativa, según que la Constitución esté o no directamente involucrada.
La clasificación es una herramienta argumentativa; lo decisivo es siempre
demostrar el derecho federal vulnerado y su relación directa e inmediata con el
resultado del pleito.} \\
\bottomrule
\end{longtable}


\chapter{Arbitrariedad de sentencia. Supuestos y alcance de la revisión}
\label{cap:p6}

\section{Una vía de revisión más allá de la cuestión federal}

Además de la cuestión federal desarrollada en el Capítulo~\ref{cap:p5}, la
jurisprudencia de la Corte Suprema construyó, sin recepción legal expresa, otra vía
por la cual es posible acudir en revisión ante ella aun cuando no se configure una
típica cuestión federal.

\begin{defbox}[: Doctrina de la arbitrariedad de sentencia]
Doctrina de creación pretoriana, desarrollada por la Corte Suprema desde el siglo
XIX, mediante la cual se atribuye a la Corte un carácter revisor de las sentencias
de tribunales inferiores aun cuando no estén configurados los presupuestos de una
cuestión federal típica, en aquellos casos en que las conclusiones de la sentencia
impugnada resultan contrarias a la sana crítica y al principio de razonabilidad.
Se engloba bajo este nombre un conjunto de supuestos que deben entenderse siempre
como una ruptura de la cadena racional y lógica de la actividad jurisdiccional.
\end{defbox}

A diferencia de la cuestión federal —regulada por la Ley 48—, la arbitrariedad de
sentencia no tiene una recepción legal específica: funciona como una válvula de
escape de creación jurisprudencial. La Corte interviene cuando detecta un
``ruido'' en la racionalidad de la decisión judicial, recurriendo con frecuencia
también al concepto de razonabilidad. Puede presentarse incluso en cuestiones de
derecho común, ajenas en principio a la competencia de la Corte, funcionando
entonces como un salvoconducto excepcional para lograr su revisión.

\section{Supuestos de arbitrariedad}

\subsection{Errores en la valoración de la prueba}

\begin{exbox}
Un testigo situado en tiempo y espacio en el lugar de los hechos, calificado y sin
circunstancias que generen demérito de su declaración, es deliberadamente omitido
en la decisión judicial. Ese tipo de omisión configura un presupuesto de
arbitrariedad. Otros supuestos análogos: una prueba pericial no debidamente
merituada, o la consideración exclusiva de una prueba que adolece de errores
graves.
\end{exbox}

\subsection{Errores en el encuadramiento jurídico}

\begin{table}[H]
\centering
\caption{Supuestos de incorrecto encuadramiento jurídico}
\begin{tabularx}{\textwidth}{@{}lX@{}}
\toprule
\textbf{Supuesto} & \textbf{Descripción} \\
\midrule
Marco jurídico inaplicable & Prescindir, por ejemplo, de los principios de la ley
de defensa del consumidor en una relación de consumo \\
Ley no vigente & Aplicar una ley antes de su entrada en vigencia, o una ley ya
derogada \\
Ausencia de fundamentación & Brindar una fundamentación totalmente ajena a la
resolución de la litis \\
\bottomrule
\end{tabularx}
\end{table}

\subsection{Violación del principio de congruencia}

\begin{exbox}
\textbf{Resolución en más}: se demanda por daños y perjuicios reclamando lucro
cesante y daño emergente, y la sentencia agrega además la pérdida de chance, no
solicitada. \textbf{Resolución en menos}: se solicita la nulidad de un acto
administrativo y una indemnización; la sentencia hace lugar a la nulidad, pero no
se adentra en la indemnización pedida. \textbf{Legitimación indebida}: se solicita
la citación de un tercero solo con ese carácter, y la sentencia resuelve una
condena en su contra.
\end{exbox}

\subsection{Exceso del tribunal de Cámara respecto de los agravios}

\begin{impbox}
Cuando existe doble instancia, la instancia revisora debe estar estrictamente
delimitada por el ámbito de los agravios expresados. Todo lo que no fue objeto de
una expresión de agravios concreta precluye, es decir, queda firme. Si la Cámara
trata elementos que no fueron objeto de impugnación, rompe el orden lógico de su
competencia y da lugar a la arbitrariedad de sentencia.
\end{impbox}

\section{Alcance de la revisión: honorarios, intereses y coexistencia con la cuestión federal}

Las regulaciones de honorarios y de intereses no motivan, en principio, la
intervención de la Corte, porque no está en juego una cuestión federal; pero si
están fijados de manera irrazonable —violando la norma procesal aplicable, o por
una capitalización derivada de un incorrecto encuadramiento jurídico—, la Corte
puede intervenir por la vía de la arbitrariedad, y de hecho lo ha hecho en
numerosos casos en materia laboral.

Cuestión federal y arbitrariedad de sentencia no son compartimentos excluyentes:
pueden convivir en una misma apelación, con lo cual el campo de revisión de la
Corte resulta, en los hechos, todavía más amplio que si solo existiera la vía de
la cuestión federal.

\section{Cuadro Cornell: arbitrariedad de sentencia}

\begin{longtable}{@{}p{0.38\textwidth}p{0.55\textwidth}@{}}
\toprule
\textbf{Preguntas / palabras clave} & \textbf{Notas / respuestas} \\
\midrule
\endhead
¿Qué es la doctrina de la arbitrariedad de sentencia? & Una doctrina de creación
jurisprudencial —no regulada por ley, a diferencia de la cuestión federal— por la
cual la Corte revisa sentencias inferiores cuando hay una ruptura de la cadena
racional y lógica de la decisión, aun sin cuestión federal. \\
¿Qué errores en la valoración de la prueba habilitan la arbitrariedad? & La
omisión deliberada de una prueba pertinente y clave, una prueba pericial no
merituada o la consideración exclusiva de una prueba con errores graves. \\
¿Cómo se viola el principio de congruencia? & Resolviendo más de lo pedido, menos
de lo pedido, o condenando a un tercero citado solo con ese carácter. \\
¿Puede la Corte intervenir en honorarios e intereses? & Sí, excepcionalmente,
cuando están fijados de forma irrazonable, aunque a priori no susciten cuestión
federal. \\
¿Pueden coexistir cuestión federal y arbitrariedad? & Sí, en una misma apelación,
ampliando el campo de revisión de la Corte. \\
\midrule
\multicolumn{2}{p{0.95\textwidth}}{\textbf{Resumen:} la arbitrariedad de
sentencia es una válvula de escape jurisprudencial que permite a la Corte revisar
sentencias de derecho común sin cuestión federal, cuando se rompe la cadena
racional de la decisión: errores graves de prueba, incorrecto encuadramiento
jurídico, falta de fundamentación, incongruencia o exceso de la Cámara sobre los
agravios. Puede coexistir con una cuestión federal y debe invocarse señalando con
precisión el error.} \\
\bottomrule
\end{longtable}


\chapter{Control de constitucionalidad de oficio. Fundamentos y precedentes}
\label{cap:p7}

\section{Concepto y fundamento}

La Corte Suprema ha aceptado la declaración de inconstitucionalidad de oficio, es
decir, aun cuando ninguna de las partes haya invocado expresamente la invalidez
constitucional de una norma. Se la concibe como un deber del órgano jurisdiccional,
no como una mera facultad discrecional.

\begin{defbox}[: \textit{Iura novit curia}]
Principio según el cual ``el juez conoce el derecho''. Habilita al tribunal a
aplicar las normas jurídicas que correspondan al caso —incluida su eventual
inconstitucionalidad— aun cuando las partes no las hayan invocado expresamente,
porque calificar jurídicamente los hechos alegados es una función propia del
juez, no de las partes.
\end{defbox}

A partir del precedente ``Mill de Pereyra, Rita Aurora y otros c/ Provincia de
Tucumán'' (CSJN, 2001), la Corte sostuvo que los jueces pueden declarar de oficio
la inconstitucionalidad de una norma, aun cuando ninguna de las partes lo haya
solicitado expresamente, sin que ello implique una violación del derecho de
defensa en juicio ni un exceso de las facultades jurisdiccionales. La Corte no
limitó esta posibilidad a determinadas materias: la sostuvo como un principio
general, aplicable a cualquier tipo de proceso.

Los fundamentos de esta facultad son tres: el principio \textit{iura novit
curia}; la competencia y el deber judicial de aplicar el derecho y de proteger la
supremacía de la Constitución, aun sin petición expresa de parte; y la conclusión
de que, por ese solo hecho, no hay violación del derecho de defensa de las
partes, porque estas ya tuvieron oportunidad de discutir los hechos y la
pretensión de fondo.

\section{El límite: la congruencia con la pretensión de la parte}

Existe, sin embargo, un límite fundamental: la eventual declaración de invalidez
debe ser coincidente con la pretensión jurídica de la parte actora, tanto al
demandar como al recurrir. El tribunal puede suplir la invocación del derecho
aplicable, pero no puede resolver más de lo que fue efectivamente pedido.

\begin{exbox}
Un trabajador demanda una indemnización por despido conforme a los topes legales
vigentes. ¿Puede el juez declarar de oficio la inconstitucionalidad de esos topes
y otorgar una suma mayor a la reclamada? La respuesta es no, porque el juez
estaría yendo más allá de lo pedido, lo que sí constituye una violación del debido
proceso: no por declarar de oficio la invalidez de la norma, sino por conceder más
de lo solicitado por la parte.
\end{exbox}

\begin{impbox}
La declaración de invalidez de una norma puede realizarse de oficio, pero siempre
que sea coincidente con la pretensión jurídica de la parte actora. Si el juez va
más allá de lo que la parte pidió, es allí —y no en la sola declaración de
oficio— donde puede configurarse una violación del derecho de defensa.
\end{impbox}

\section{Precedentes históricos}

\textit{Marbury v.\ Madison} y \textit{Sojo} (desarrollados en el
Capítulo~\ref{cap:p2}) son, en cierto sentido, antecedentes remotos de esta
doctrina: en ambos casos la invalidez —o, en \textit{Sojo}, la cuestión de
competencia— fue resuelta por el tribunal sin que la parte hubiera argumentado
estrictamente en esos términos, mostrando un camino largo y relativamente pacífico
hacia la doctrina actual. Otro precedente relevante es el caso de la
``Municipalidad de la Capital c/ Elortondo'', referido a la expropiación,
considerado uno de los primeros fallos en que la Corte se inmiscuye en la
regulación económica de contratos de derecho privado, también bajo la condición
de que la invalidación resultara congruente con la pretensión jurídica de alguna
de las partes.

\section{Vínculo con el recurso extraordinario federal}

La facultad de declarar la inconstitucionalidad de oficio, si bien se estudia
como un tema autónomo, con frecuencia se activa en el propio trámite del recurso
extraordinario: el juez, al conocer el derecho, puede advertir una colisión
normativa no planteada expresamente por las partes, siempre que la declaración
resulte coincidente con la pretensión jurídica deducida en el proceso.

\section{Cuadro Cornell: inconstitucionalidad de oficio}

\begin{longtable}{@{}p{0.38\textwidth}p{0.55\textwidth}@{}}
\toprule
\textbf{Preguntas / palabras clave} & \textbf{Notas / respuestas} \\
\midrule
\endhead
¿Qué es la declaración de inconstitucionalidad de oficio? & La facultad,
concebida como deber del órgano jurisdiccional, de declarar la invalidez
constitucional de una norma aun cuando ninguna de las partes lo haya solicitado
expresamente, a partir de ``Mill de Pereyra'' (2001). \\
¿En qué principio se funda y por qué no viola el derecho de defensa? & En el
principio \textit{iura novit curia}: el juez conoce el derecho. Al tratarse de
una calificación jurídica y no de nuevos hechos, en principio no hay violación
del derecho de defensa. \\
¿Cuál es el límite de esta facultad? & La pretensión de la parte: el tribunal
puede suplir la invocación del derecho, pero no puede resolver por encima de lo
que la parte solicitó. \\
¿Qué precedentes se vinculan con esta doctrina? & \textit{Marbury v.\ Madison} y
\textit{Sojo}, donde la invalidez se resolvió sin haber sido invocada en esos
términos exactos; y el caso ``Elortondo'', sobre expropiación. \\
\midrule
\multicolumn{2}{p{0.95\textwidth}}{\textbf{Resumen:} los tribunales pueden
declarar de oficio la inconstitucionalidad de una norma en virtud del principio
\textit{iura novit curia}, sin que ello viole el derecho de defensa, siempre que
la declaración sea congruente con la pretensión de la parte al demandar y al
recurrir.} \\
\bottomrule
\end{longtable}

\chapter{Recurso extraordinario: requisitos formales, comunes y propios}
\label{cap:p8}

\section{El recurso extraordinario como mecanismo de impugnación}

El recurso extraordinario federal es el mecanismo procesal, regulado por la Ley
48, por el cual una cuestión constitucional resuelta por un tribunal superior de
una causa —provincial o nacional— puede llegar en revisión hasta la Corte Suprema
de Justicia de la Nación. Es la vía de acceso ordinaria a la Corte cuando existe
una cuestión federal (Capítulo~\ref{cap:p5}).

\subsection{Requisitos formales}

La Corte reglamentó minuciosamente la forma del recurso mediante la
\textbf{Acordada 4/2007}, que fija reglas estrictas de admisibilidad formal: una
extensión máxima de 40 páginas, tipografía y espaciado determinados, una carátula
con un resumen del caso y de la cuestión federal invocada, entre otros recaudos.
El incumplimiento de estos requisitos formales puede llevar al rechazo del recurso
sin que se examine su fondo.

El recurso debe interponerse ante el mismo tribunal que dictó la sentencia
impugnada (el ``tribunal superior de la causa''), dentro del plazo de diez días
desde la notificación de la sentencia. Un requisito adicional, exigido por la
jurisprudencia, es el de la \textbf{reserva}: la cuestión federal debe haberse
planteado ya en cada instancia anterior del proceso, en la primera oportunidad
procesal posible; no puede introducirse por primera vez recién al interponer el
recurso extraordinario.

\subsection{Requisitos comunes a todo recurso judicial}

\begin{defbox}[: Gravamen]
Perjuicio acreditado en cabeza de quien impugna. La decisión cuestionada debe
generar, inevitablemente, un perjuicio concreto y actual a quien recurre; no puede
ser conjetural ni referido a un daño meramente futuro y eventual.
\end{defbox}

Debe existir, además, un \textbf{proceso judicial en trámite} —en contraposición,
por ejemplo, a una instancia meramente administrativa—, presupuesto ineludible
para la procedencia de cualquier recurso, incluido el extraordinario.

\subsection{Requisitos propios: sentencia definitiva o asimilable}

\begin{defbox}[: Sentencia definitiva o asimilable]
El recurso extraordinario, en principio, no está previsto para actos
interlocutorios ni para medidas cautelares. Debe tratarse de sentencias que ponen
fin al proceso, o asimilables a definitivas: aquellas en las que ya no hay marcha
atrás, o solo hay marcha hacia adelante, generando un perjuicio irreparable si no
se revisan de inmediato.
\end{defbox}

Las medidas cautelares (desarrolladas en profundidad en el Capítulo~\ref{cap:p10})
son, por definición, provisionales: no son sentencias, porque cualquier medida
cautelar debe luego ser ratificada por la decisión de fondo. Por eso, en
principio, el recurso extraordinario no procede contra ellas. Excepcionalmente, en
casos de especial trascendencia, gravedad institucional o que involucran una
pluralidad de sujetos con fuerte impacto social, la Corte se aboca a recursos
extraordinarios interpuestos contra medidas cautelares, cuando estas pueden
generar perjuicios irreparables.

\begin{exbox}
Cuando se sancionó la Ley de Servicios de Comunicación Audiovisual (``Ley de
Medios''), el grupo empresario afectado la impugnó e interpuso una medida cautelar
para suspender su aplicación. La Corte abordó el recurso extraordinario pese a
estar dirigido contra una cautelar, por tratarse de una ley recientemente dictada
con un enorme efecto social e institucional.
\end{exbox}

\begin{table}[H]
\centering
\caption{Síntesis de los requisitos del recurso extraordinario}
\begin{tabularx}{\textwidth}{@{}lX@{}}
\toprule
\textbf{Tipo de requisito} & \textbf{Contenido} \\
\midrule
Formales & Acordada 4/2007: extensión máxima de 40 páginas, forma de la carátula
y del escrito; plazo de 10 días; reserva del caso federal en cada instancia
anterior \\
Comunes & Gravamen (perjuicio actual y concreto) y existencia de un proceso
judicial en trámite \\
Propios & Sentencia definitiva o asimilable; existencia de una cuestión federal
(Capítulo~\ref{cap:p5}), con relación directa e inmediata y decisión contraria al
derecho federal invocado \\
\bottomrule
\end{tabularx}
\end{table}

\begin{impbox}
Efectos de la concesión y del rechazo. Si el tribunal concede el recurso, la
concesión tiene, en principio, efecto suspensivo sobre la ejecución de la
sentencia. Si el tribunal lo rechaza, no hay efecto suspensivo: la sentencia queda
en condiciones de ejecutarse, sin perjuicio de que la parte recurra ese rechazo
mediante la queja (Capítulo~\ref{cap:p9}).
\end{impbox}

\section{Cuadro Cornell: requisitos del recurso extraordinario}

\begin{longtable}{@{}p{0.38\textwidth}p{0.55\textwidth}@{}}
\toprule
\textbf{Preguntas / palabras clave} & \textbf{Notas / respuestas} \\
\midrule
\endhead
¿Qué regula la Acordada 4/2007? & Los requisitos formales del recurso
extraordinario: extensión máxima, forma de la carátula y del escrito. \\
¿Qué exigen los requisitos comunes a todo recurso? & Un gravamen o perjuicio
actual y concreto, y la existencia de un proceso judicial en trámite. \\
¿Qué tipo de resoluciones habilita el recurso extraordinario? & Sentencias
definitivas o asimilables a definitivas; en principio no medidas cautelares,
salvo gravedad institucional o fuerte impacto social. \\
¿Qué es la ``reserva'' del caso federal? & La exigencia de haber planteado la
cuestión federal en cada instancia anterior, en la primera oportunidad procesal
posible. \\
¿Qué efecto tiene la concesión del recurso? & Efecto suspensivo, en principio,
sobre la ejecución de la sentencia impugnada. \\
\midrule
\multicolumn{2}{p{0.95\textwidth}}{\textbf{Resumen:} el recurso extraordinario
exige requisitos formales (Acordada 4/2007, plazo de diez días, reserva del caso
federal), requisitos comunes (gravamen y proceso judicial en trámite) y
requisitos propios (sentencia definitiva o asimilable y cuestión federal, con
relación directa e inmediata). Su concesión suspende, en principio, la ejecución
de la sentencia impugnada.} \\
\bottomrule
\end{longtable}


\chapter{Recurso de queja}
\label{cap:p9}

\section{Cuándo procede}

\begin{defbox}[: Recurso de queja por recurso extraordinario denegado]
Recurso que interpone la parte que dedujo el recurso extraordinario, si el
tribunal superior de la causa lo rechaza y ella quiere insistir en su
tratamiento. Se interpone directamente ante la Corte Suprema, sin pasar por el
tribunal que dictó el rechazo.
\end{defbox}

Mientras que el recurso extraordinario no está sujeto al pago de una tasa, la
queja sí requiere un depósito previo (salvo excepciones, como la del trabajador en
procesos laborales); si la queja no es admitida, ese monto se destina a la
biblioteca de la Corte. El plazo para interponerla es de cinco días desde la
notificación del rechazo, ampliable según la tabla de distancias para quien tiene
su domicilio en el interior del país, en un escrito de hasta diez páginas.

\begin{table}[H]
\centering
\caption{Recurso extraordinario y recurso de queja}
\begin{tabularx}{\textwidth}{@{}lXX@{}}
\toprule
& \textbf{Recurso extraordinario} & \textbf{Queja} \\
\midrule
Ante quién & El mismo tribunal que dictó el acto impugnado & Directamente ante la
Corte Suprema \\
Costo & No se cobra tasa & Requiere depósito, salvo excepciones \\
Extensión máxima & 40 páginas & 10 páginas \\
Plazo & 10 días & 5 días (más tabla de distancias) \\
Cuándo procede & Contra sentencia definitiva o asimilable & Cuando el recurso
extraordinario fue rechazado \\
\bottomrule
\end{tabularx}
\end{table}

\section{Objeto de la queja y efectos}

El centro de la argumentación de la queja consiste en demostrar que el rechazo del
recurso extraordinario fue arbitrario: infundado, ilegítimo o irrazonable. A
diferencia del recurso extraordinario concedido, la mera interposición de la
queja no suspende la ejecución de la sentencia, porque la Corte la resuelve
\textit{inaudita parte} (sin dar traslado a la contraparte). Para intentar
suspender la ejecución mientras la queja tramita, la única vía disponible es
solicitar una medida cautelar, sustentada en la verosimilitud del derecho a
interponer el recurso extraordinario, tema que se desarrolla en el
Capítulo~\ref{cap:p10}.

\begin{impbox}
Relación con el control de constitucionalidad de oficio. Como se explicó en el
Capítulo~\ref{cap:p7}, la facultad de declarar la inconstitucionalidad de oficio
también puede activarse en el marco de una queja: la Corte, al resolverla, puede
advertir una colisión normativa no planteada expresamente, siempre que la
declaración sea congruente con la pretensión deducida en el proceso.
\end{impbox}

\section{Cuadro Cornell: recurso de queja}

\begin{longtable}{@{}p{0.38\textwidth}p{0.55\textwidth}@{}}
\toprule
\textbf{Preguntas / palabras clave} & \textbf{Notas / respuestas} \\
\midrule
\endhead
¿Qué es el recurso de queja? & El recurso que se interpone directamente ante la
Corte Suprema cuando el tribunal superior de la causa rechazó el recurso
extraordinario. \\
¿Qué plazo y extensión tiene? & Cinco días desde la notificación del rechazo (más
tabla de distancias), en un escrito de hasta diez páginas. \\
¿Requiere depósito? & Sí, salvo excepciones como la del trabajador; si se
rechaza, el monto se destina a la biblioteca de la Corte. \\
¿Suspende la ejecución de la sentencia? & No, a diferencia del recurso
extraordinario concedido; la Corte resuelve \textit{inaudita parte}. \\
¿Cómo se puede suspender la ejecución mientras tramita la queja? & Solicitando
una medida cautelar fundada en la verosimilitud del derecho a interponer el
recurso extraordinario. \\
\midrule
\multicolumn{2}{p{0.95\textwidth}}{\textbf{Resumen:} la queja es la vía para
insistir ante la propia Corte Suprema cuando el recurso extraordinario fue
rechazado por el tribunal inferior. Exige depósito previo, un plazo breve de
cinco días y un escrito de hasta diez páginas, y su objeto es demostrar que el
rechazo fue arbitrario. A diferencia del recurso extraordinario concedido, no
suspende por sí sola la ejecución de la sentencia.} \\
\bottomrule
\end{longtable}


\chapter{Medida cautelar. Presupuestos. Su relación con el recurso extraordinario}
\label{cap:p10}

\section{Concepto y función}

\begin{defbox}[: Medida cautelar]
Medida provisional, con un límite temporal, que no es una decisión definitiva ni
hace cosa juzgada. Puede solicitarse en cualquier momento del proceso, incluso
antes de su inicio. Su función es dar validez, ejecutoriedad y virtualidad a la
eventual decisión sobre el fondo del asunto, evitando que el mero transcurso del
tiempo genere perjuicios irreparables mientras el proceso continúa.
\end{defbox}

\section{Los tres presupuestos clásicos}

\begin{enumerate}[leftmargin=1.6em]
  \item \textbf{Verosimilitud en el derecho}: que el planteo tenga asidero
  jurídico razonable; no se exige una certeza absoluta, sino un ``humo del buen
  derecho'' suficiente para habilitar una decisión provisoria.
  \item \textbf{Peligro en la demora}: la demostración de que el transcurso del
  tiempo puede hacer perder un derecho, generando un perjuicio irreparable o de
  muy dificultosa reparación posterior.
  \item \textbf{Contracautela}: la garantía que debe prestar quien solicita la
  medida, para responder por los eventuales daños que su dictado pueda causar a
  la contraparte si finalmente la pretensión de fondo no prospera.
\end{enumerate}

\begin{exbox}
Si una persona necesita acceder a una prestación médica de alta complejidad,
puede solicitar cautelarmente la prestación mientras se discute si estaba a cargo
de la obra social o de la empresa de medicina prepaga: de lo contrario, el mero
transcurso del tiempo del proceso puede generar un daño irreparable en su salud.
\end{exbox}

\begin{impbox}
Vasos comunicantes. Verosimilitud en el derecho y peligro en la demora se
relacionan de forma inversamente proporcional: cuanto mayor es la verosimilitud
del derecho invocado, menos exigente puede ser la demostración del peligro en la
demora, y viceversa. Ninguno de los dos requisitos puede faltar por completo: el
tribunal los pondera siempre en conjunto.
\end{impbox}

\section{Relación con el recurso extraordinario}

Como se adelantó en el Capítulo~\ref{cap:p8}, el recurso extraordinario no procede,
en principio, contra resoluciones cautelares, precisamente porque no son
sentencias definitivas. Sin embargo, cuando la cautelar tiene efectos generales,
excede el interés puntual de las partes y puede configurar un agravio irreparable,
muy excepcionalmente se habilita el recurso extraordinario en su contra —tal como
ocurrió en el caso de la Ley de Medios, o en controversias sobre el financiamiento
de la universidad pública, donde el fuero contencioso administrativo hizo lugar a
una cautelar pero habilitó su revisión por la Corte dada la delicadeza política del
caso—.

Contra el Estado Nacional rige, además, un régimen especial: la Ley 26.854 prevé
un traslado previo, urgente y por un plazo breve, a la autoridad pública
demandada, antes de resolver sobre la cautelar, a diferencia de la regla general
de resolución \textit{inaudita parte} que rige para la queja (Capítulo~\ref{cap:p9}).

\section{Cuadro Cornell: medida cautelar}

\begin{longtable}{@{}p{0.38\textwidth}p{0.55\textwidth}@{}}
\toprule
\textbf{Preguntas / palabras clave} & \textbf{Notas / respuestas} \\
\midrule
\endhead
¿Qué es una medida cautelar? & Una medida provisional, sin efecto de cosa
juzgada, destinada a asegurar que la eventual decisión de fondo pueda cumplirse
efectivamente. \\
¿Cuáles son sus tres presupuestos clásicos? & Verosimilitud en el derecho,
peligro en la demora y contracautela. \\
¿Cómo se relacionan verosimilitud y peligro en la demora? & De forma inversamente
proporcional: a mayor verosimilitud, menor exigencia de peligro en la demora, y
viceversa. \\
¿Procede el recurso extraordinario contra una cautelar? & En principio no, por no
ser sentencia definitiva; excepcionalmente sí, cuando hay gravedad institucional o
efectos generales con perjuicio irreparable. \\
¿Qué régimen especial rige contra el Estado Nacional? & La Ley 26.854, que exige
un traslado previo y urgente a la autoridad pública antes de resolver la cautelar. \\
\midrule
\multicolumn{2}{p{0.95\textwidth}}{\textbf{Resumen:} la medida cautelar asegura
la eficacia práctica del proceso mientras se resuelve el fondo, sobre la base de
verosimilitud del derecho, peligro en la demora y contracautela. Aunque el
recurso extraordinario no procede, en principio, contra ella por no ser sentencia
definitiva, la gravedad institucional o el impacto general de la medida pueden
habilitar excepcionalmente su revisión por la Corte.} \\
\bottomrule
\end{longtable}


\chapter{\textit{Per saltum}}
\label{cap:p11}

\section{Concepto y marco normativo}

\begin{defbox}[: \textit{Per saltum}]
Figura prevista en los artículos 257 bis y 257 ter del Código Procesal Civil y
Comercial de la Nación, que permite, en el ámbito federal, pasar directamente de
la primera instancia a la Corte Suprema, sin pasar por la cámara de apelaciones,
cuando la cuestión debatida sea de notoria gravedad institucional cuya solución
definitiva y expedita sea necesaria, y el recurso constituya el único remedio
eficaz para la protección del derecho federal comprometido.
\end{defbox}

\section{Requisitos}

\begin{enumerate}[leftmargin=1.6em]
  \item Causa de índole federal.
  \item Gravedad institucional que exceda el interés particular de las partes:
  el caso, en general, debe proyectarse sobre una controversia de enorme
  relevancia social, porque atañe a muchos sujetos que pueden estar en la misma
  situación.
  \item Riesgo de un perjuicio irreparable si el caso sigue el trámite recursivo
  ordinario, es decir, si debe pasar primero por la cámara de apelaciones.
\end{enumerate}

\begin{exbox}
Un juez de primera instancia había declarado la inconstitucionalidad de un
artículo de una reforma laboral; el gobierno solicitó el \textit{per saltum}
frente a la medida cautelar que había suspendido varios artículos de la ley. La
Corte rechazó la pretensión al no encontrar acreditada la extrema celeridad
invocada, y consideró razonable que la causa siguiera su curso ordinario ante la
cámara, que finalmente levantó la suspensión. En un caso similar, referido a una
ley de empleo público, la Corte también rechazó el \textit{per saltum} por no
encontrar acreditada la gravedad institucional ni el perjuicio irreparable,
reafirmando el carácter absolutamente excepcional de esta figura.
\end{exbox}

\begin{impbox}
El \textit{per saltum} está previsto solo para causas de derecho federal: nunca
está en discusión que se salte la instancia de una cámara provincial hacia la
Corte provincial, porque todo lo que es materia federal tramita, desde el inicio,
ante la justicia federal.
\end{impbox}

\section{Cuadro Cornell: \textit{per saltum}}

\begin{longtable}{@{}p{0.38\textwidth}p{0.55\textwidth}@{}}
\toprule
\textbf{Preguntas / palabras clave} & \textbf{Notas / respuestas} \\
\midrule
\endhead
¿Qué es el \textit{per saltum}? & Una figura excepcional (arts.\ 257 bis y ter
CPCCN) que permite saltear la instancia de la cámara federal y acceder
directamente a la Corte, en causas de derecho federal con gravedad institucional
y riesgo de perjuicio irreparable. \\
¿Cuáles son sus tres requisitos? & Causa de índole federal, gravedad
institucional que exceda a las partes, y riesgo de perjuicio irreparable por el
trámite ordinario. \\
¿Se aplica en causas provinciales? & No; está previsto solo para procesos de
derecho federal. \\
¿Cómo lo trató la Corte en los casos de reforma laboral y empleo público? & Los
rechazó en ambos casos, por no encontrar acreditada la extrema gravedad
institucional ni el perjuicio irreparable, reafirmando su carácter excepcional. \\
\midrule
\multicolumn{2}{p{0.95\textwidth}}{\textbf{Resumen:} el \textit{per saltum} es
una vía excepcional para llegar directamente a la Corte Suprema sin pasar por la
cámara, reservada a causas federales de notoria gravedad institucional con
riesgo de perjuicio irreparable. La jurisprudencia lo aplica con suma cautela,
como muestran los rechazos en materia de reforma laboral y de empleo público.} \\
\bottomrule
\end{longtable}

% ============================================================
% PARTE IV - EL PRINCIPIO DE RESERVA Y LOS DERECHOS INDIVIDUALES
% ============================================================
\part{El principio de reserva y los derechos individuales}

\chapter{Principio de reserva. Objeción de conciencia}
\label{cap:p12}

\section{El artículo 19: el principio de reserva}

\begin{defbox}[: Artículo 19 de la Constitución Nacional]
``Las acciones privadas de los hombres que de ningún modo ofendan al orden y a la
moral pública, ni perjudiquen a un tercero, están sólo reservadas a Dios, y
exentas de la autoridad de los magistrados. Ningún habitante de la Nación será
obligado a hacer lo que no manda la ley, ni privado de lo que ella no prohíbe.''
\end{defbox}

\begin{defbox}[: Principio de reserva]
Ámbito de la conducta humana vedado a la regulación normativa: las acciones
privadas que no ofenden el orden ni la moral pública ni perjudican a terceros. El
Congreso no puede regular conductas que no exteriorizan una afectación a terceros.
El punto más complejo de esta garantía consiste en determinar dónde está la veda
normativa y cuál es el verdadero umbral en el que las normas no pueden
entrometerse.
\end{defbox}

Es, en palabras de Kelsen, una típica \textbf{norma de cerramiento}: brinda un
principio interpretativo para entender dónde están las fronteras de la regulación
legal, pero no resuelve por sí sola los casos particulares. La Constitución da una
idea general, y es la interpretación judicial —caso a caso— la que traza,
efectivamente, esa frontera. Entre los campos de aplicación de este principio se
destacan la no punición del consumo personal de estupefacientes (doctrina
\textit{Bazterrica}, CSJN, 1986), la interrupción del embarazo, la muerte asistida
y, en lo que aquí interesa especialmente, la \textbf{libertad de conciencia y de
religión}, desarrollada en el Capítulo~\ref{cap:p13}.

\begin{impbox}
Relación con la interpretación constitucional. El propio artículo 19 es un buen
ejemplo del problema estudiado en el Capítulo~\ref{cap:p1}: una lectura
textualista tiende a acotar la interpretación a la literalidad de ``orden'',
``moral pública'' y ``perjuicio a terceros''; una lectura no textualista (sea
originalista o pragmática) admite proyectar esas expresiones a los valores y
problemas actuales.
\end{impbox}

\section{La objeción de conciencia}

\begin{notebox}[Ampliación externa]
Los apuntes de la cátedra no traían un desarrollo autónomo de la objeción de
conciencia como categoría general, aunque mencionaban de modo tangencial la
objeción religiosa al servicio militar dentro del capítulo sobre libertad
religiosa. El desarrollo que sigue —incluidos los casos ``Portillo'' y
``Bahamondez''— se elaboró a partir de fuentes jurisprudenciales y doctrinarias
verificadas (véase la bibliografía final). No fue posible verificar con las
fuentes disponibles el contenido ni la carátula exacta de un tercer precedente
mencionado en el cronograma de la cátedra bajo el apellido ``Sisto''; por
prudencia académica, se deja constancia de esa laguna en lugar de completarla con
información no verificada.
\end{notebox}

\begin{defbox}[: Objeción de conciencia]
Derecho de una persona a ser eximida del cumplimiento de un deber legal —o a
cumplirlo por una vía alternativa compatible con sus convicciones— cuando ese
deber contradice creencias religiosas o morales profundas, serias y acreditadas.
Se trata de una tensión entre un derecho constitucional (la libertad de culto y
de conciencia) y un deber también de raigambre constitucional o legal, que debe
resolverse armonizando ambos extremos y no haciendo prevalecer automáticamente a
uno sobre el otro.
\end{defbox}

La objeción de conciencia se ubica dentro del principio de reserva del artículo
19, en tanto protege un ámbito de convicciones íntimas de la persona, pero se
proyecta hacia afuera cuando esas convicciones chocan con un deber jurídico
positivo (por ejemplo, el deber de armarse en defensa de la Nación, art.\ 21 CN,
o el deber de someterse a un tratamiento médico). Por eso su reconocimiento
constitucional deriva, en la jurisprudencia argentina, de la conjunción entre el
artículo 19 (autonomía personal) y el artículo 14 (libertad de culto y de
conciencia) de la Constitución Nacional.

\subsection{El caso ``Portillo'' (CSJN, 1989)}

\begin{exbox}
Alfredo Portillo, de religión católica, fue condenado por no presentarse a
cumplir el servicio militar obligatorio (art.\ 44, Ley 17.531), alegando que sus
convicciones religiosas —fundadas en el quinto mandamiento— le impedían portar
armas contra otro ser humano. La Corte Suprema, por mayoría (jueces Fayt,
Petracchi y Bacqué), reconoció por primera vez en la Argentina el carácter
constitucional de la objeción de conciencia, derivándola de la libertad de culto y
de conciencia del artículo 14 CN. Sin embargo, no lo eximió lisa y llanamente del
servicio: armonizó ese derecho con el deber de armarse en defensa de la Patria y
de la Constitución (art.\ 21 CN), confirmando la obligación de Portillo de
cumplir el servicio militar, pero \textbf{sin el uso de armas}, en una modalidad
compatible con sus convicciones religiosas.
\end{exbox}

\begin{impbox}
La clave de ``Portillo'' es la armonización, no la exención automática. Ni los
derechos ni los deberes constitucionales son absolutos: obligar a quien profesa
una objeción religiosa seria a portar armas lesiona su conciencia, pero eximirlo
lisa y llanamente de toda carga pública, sin más, rompería la igualdad en las
cargas públicas frente a quienes sí cumplen el servicio en armas. La solución de
la Corte fue, por eso, una vía intermedia: mantener la carga, pero adaptar su
modo de cumplimiento.
\end{impbox}

\subsection{El caso ``Bahamondez'' (CSJN, 1993)}

\begin{exbox}
Marcelo Bahamondez, internado por una hemorragia digestiva, se negó a recibir
transfusiones de sangre por ser Testigo de Jehová, religión que las considera
pecaminosas. Las instancias inferiores autorizaron la transfusión en contra de su
voluntad, calificando su negativa como un ``suicidio lentificado''. Al llegar el
caso a la Corte Suprema, la situación de hecho ya había cesado (Bahamondez había
recibido el alta médica sin haber sido transfundido), por lo que la mayoría
declaró la cuestión abstracta. Sin embargo, distintos votos —entre ellos los de
los jueces Barra, Fayt, Belluscio y Petracchi— avanzaron sobre el fondo del
asunto, sosteniendo que la negativa de un adulto, lúcido, a un tratamiento médico
por razones de conciencia religiosa, sin afectación a terceros, se encuentra
amparada por el principio de reserva del artículo 19 CN y por la autonomía
personal, y que el Estado no puede imponerle el tratamiento contra su voluntad
invocando la protección de su propia vida, salvo que existan derechos de terceros
comprometidos.
\end{exbox}

\begin{table}[H]
\centering
\caption{``Portillo'' y ``Bahamondez'': dos variantes de objeción de conciencia}
\begin{tabularx}{\textwidth}{@{}lXX@{}}
\toprule
& \textbf{Portillo (1989)} & \textbf{Bahamondez (1993)} \\
\midrule
Deber en tensión & Servicio militar obligatorio (art.\ 21 CN) & Deber estatal de
tutelar la vida y la salud \\
Fundamento de la objeción & Libertad de culto y conciencia (art.\ 14 CN),
motivación católica & Libertad de culto y autonomía personal (arts.\ 14 y 19 CN),
motivación como Testigo de Jehová \\
Solución & Armonización: se mantiene el deber, pero se modifica su modo de
cumplimiento (servicio sin armas) & La autonomía del paciente adulto y lúcido
prevalece sobre la imposición estatal de un tratamiento, sin afectación de
terceros \\
Aporte doctrinal & Primer reconocimiento constitucional expreso de la objeción de
conciencia en la Argentina & Proyección de la objeción de conciencia sobre el
derecho a rechazar tratamientos médicos, dentro del principio de reserva \\
\bottomrule
\end{tabularx}
\end{table}

\begin{exambox}
El principio de reserva del artículo 19 CN traza la frontera entre lo que el
Estado puede y no puede regular, exigiendo que exista ofensa al orden, a la moral
pública o perjuicio a un tercero para que una conducta privada pueda ser objeto de
regulación legal. La objeción de conciencia es una proyección de ese principio,
combinado con la libertad de culto y de conciencia del artículo 14 CN: protege el
derecho de una persona a no cumplir, o a cumplir de un modo alternativo, un deber
jurídico que contradice sus convicciones religiosas o morales serias y
acreditadas. La Corte Suprema la reconoció por primera vez en ``Portillo''
(1989), armonizando el deber de armarse en defensa de la Patria con la libertad
de conciencia mediante una solución intermedia (servicio sin armas), y la
proyectó en ``Bahamondez'' (1993) hacia el derecho de un adulto lúcido a rechazar
un tratamiento médico por motivos religiosos, en tanto no se afecten derechos de
terceros. En ambos casos, la objeción de conciencia no opera como una exención
automática, sino como un llamado a armonizar el derecho invocado con el deber o
interés estatal en juego.
\end{exambox}

\section{Cuadro Cornell: principio de reserva y objeción de conciencia}

\begin{longtable}{@{}p{0.38\textwidth}p{0.55\textwidth}@{}}
\toprule
\textbf{Preguntas / palabras clave} & \textbf{Notas / respuestas} \\
\midrule
\endhead
¿Qué establece el artículo 19 CN? & Que las acciones privadas que no ofenden el
orden ni la moral pública, ni perjudican a terceros, están exentas de la autoridad
de los magistrados. \\
¿Qué es una ``norma de cerramiento''? & En palabras de Kelsen, una norma que
traza un principio interpretativo sobre las fronteras de la regulación legal,
aplicable caso a caso. \\
¿Qué es la objeción de conciencia? & El derecho a ser eximido de un deber legal, o
a cumplirlo de otro modo, cuando contradice convicciones religiosas o morales
serias y acreditadas. \\
¿Qué resolvió ``Portillo'' (1989)? & Reconoció por primera vez la objeción de
conciencia, armonizando el deber de armarse (art.\ 21 CN) con la libertad de
conciencia (art.\ 14 CN): servicio militar sin uso de armas. \\
¿Qué resolvió ``Bahamondez'' (1993)? & Proyectó la objeción de conciencia sobre
el derecho de un adulto lúcido a rechazar tratamientos médicos por motivos
religiosos, sin afectación de terceros, dentro del principio de reserva. \\
\midrule
\multicolumn{2}{p{0.95\textwidth}}{\textbf{Resumen:} el artículo 19 CN traza la
frontera entre lo regulable y lo reservado a la autonomía personal. La objeción de
conciencia, construida sobre esa base junto con la libertad de culto del artículo
14 CN, no opera como exención automática de los deberes legales, sino como una
armonización entre el deber en juego y la convicción invocada, como muestran
``Portillo'' (servicio sin armas) y ``Bahamondez'' (autonomía frente al
tratamiento médico).} \\
\bottomrule
\end{longtable}


\chapter{Libertad religiosa}
\label{cap:p13}

\section{La Iglesia como persona jurídica de derecho público}

El ordenamiento argentino reconoce a la Iglesia Católica como persona jurídica de
derecho público, lo que conlleva privilegios como la inmunidad e
inembargabilidad de los bienes afectados al culto, extendida también a otras
religiones registradas ante la Secretaría de Culto. Originalmente se exigía que el
presidente profesara la religión católica, requisito derogado por la reforma
constitucional de 1994, que mantuvo, sin embargo, el estatus de la Iglesia como
persona jurídica de derecho público. A partir de esa reforma, la idea del
sostenimiento del culto católico debe conjugarse inevitablemente con una férrea
libertad religiosa y con el principio de no discriminación entre los distintos
credos.

\begin{defbox}[: Neutralidad estatal en materia religiosa]
El sostenimiento del culto católico se reduce a los aspectos institucionales de
configuración de la Iglesia como persona jurídica de derecho público, pero no
puede trasladarse a preferencias expresas en la opción religiosa de las personas.
La profesión de una religión queda relegada al ámbito de reserva del artículo 19
CN.
\end{defbox}

\section{Categorías sospechosas}

\begin{defbox}[: Categorías sospechosas]
Regulaciones que establecen diferencias basadas en rasgos como edad, nacionalidad,
religión, origen o género. La jurisprudencia de la Corte las somete a un
escrutinio más estricto, invirtiendo la presunción de legitimidad de la norma.
\end{defbox}

Las normas emanadas de competencias legítimas se presumen, en general, válidas y
aplicables. Lo que hace la categoría sospechosa es invertir esa carga: la norma se
presume, a priori, potencialmente inválida, y es el Estado quien, en el juicio,
debe demostrar su legitimidad y razonabilidad.

\begin{impbox}
Relación con el poder judicial contramayoritario. Esta técnica traduce, en el
derecho argentino, la misma lógica que Ely desarrolla a propósito de las minorías
discretas e insulares (Capítulo~\ref{cap:p4}): cuando una norma distingue por
religión u otro rasgo sensible, el grupo afectado suele carecer de peso suficiente
en el proceso político ordinario para revertir esa norma por la vía legislativa,
lo que justifica un control judicial más exigente.
\end{impbox}

\section{El caso de la educación religiosa obligatoria en Salta}

Una demanda cuestionó la educación inicial obligatoria en la provincia de Salta,
en la que se impartía, en los hechos, catequesis católica bajo la apariencia de
una asignatura de ``religión a secas''. La asociación demandante sostenía que, al
no existir en la práctica una materia alternativa efectiva para quienes se
eximían, se generaba una \textbf{discriminación indirecta} hacia los alumnos no
católicos.

\begin{impbox}
Discriminación indirecta. Lo relevante del caso no era la existencia formal de
una opción de exención, sino su ineficacia práctica: los alumnos exentos no
recibían una materia alternativa real y quedaban, en los hechos, relegados frente
a la mayoría, por lo que la sola previsión formal de una excepción no bastaba para
salvar la norma.
\end{impbox}

La Corte, por mayoría, resolvió que el sistema, tal como estaba implementado, era
violatorio de la libertad religiosa, por establecer preferencias y evidenciar una
segregación hacia los alumnos que no cursaban la clase de religión, sin ofrecerles
una alternativa educativa real. El argumento de la autonomía provincial en materia
de educación inicial, invocado por la propia provincia, no fue suficiente para
justificar la norma frente a la protección de la libertad religiosa.

\section{Otros ámbitos de aplicación}

\subsection{Testigos de Jehová y transfusiones de sangre}

Como se desarrolló en el Capítulo~\ref{cap:p12} a propósito de ``Bahamondez'', la
negativa a recibir transfusiones de sangre por motivos religiosos queda amparada
por el principio de reserva y la libertad religiosa, siempre que se trate de una
decisión personal, inequívoca y expresada en condiciones de lucidez y libre
discernimiento, sin afectación de terceros. Esa voluntad no puede ser suplida, ni
siquiera por familiares allegados, cuando la persona está en condiciones de
expresarla por sí misma.

\subsection{Educación en el hogar y comunidades menonitas}

La Constitución reconoce que todo habitante es libre de enseñar y de aprender
(art.\ 14 CN), sin exigir que ello ocurra necesariamente en una institución
formal. Comunidades religiosas como la menonita han planteado, en distintos
países, su derecho a educar a sus hijos según sus propias pautas culturales y
religiosas, sin enviarlos a los sistemas de educación básica estatales, siempre
que se acredite el cumplimiento de contenidos mínimos obligatorios.

\subsection{Objeción religiosa al servicio militar}

Este ámbito se desarrolla en profundidad en el Capítulo~\ref{cap:p12} a propósito
del caso ``Portillo'': religiones que impiden portar armas plantearon,
históricamente, tensiones con el servicio militar obligatorio, en un contexto en
que ese servicio no era únicamente una prestación de defensa, sino que conllevaba
también funciones de sanidad pública.

\begin{defbox}[: El principio de reserva como hilo conductor]
El principio de reserva (art.\ 19 CN) constituye el eje común que atraviesa todos
estos casos: transfusiones de sangre, educación en el hogar, comunidades
menonitas y objeción religiosa al servicio militar. En todos ellos, la pregunta de
fondo es la misma que en el Capítulo~\ref{cap:p12}: dónde termina la autonomía
personal y dónde empieza el ámbito legítimo de la regulación estatal.
\end{defbox}

\section{Cuadro Cornell: libertad religiosa}

\begin{longtable}{@{}p{0.38\textwidth}p{0.55\textwidth}@{}}
\toprule
\textbf{Preguntas / palabras clave} & \textbf{Notas / respuestas} \\
\midrule
\endhead
¿Qué privilegios conlleva el estatus de la Iglesia Católica como persona jurídica
de derecho público? & La inmunidad e inembargabilidad de los bienes afectados al
culto, extendida a otras religiones registradas ante la Secretaría de Culto. \\
¿Qué es una categoría sospechosa y qué efecto produce? & Una distinción normativa
basada en rasgos como religión, nacionalidad, edad, origen o género, que invierte
la presunción de legitimidad: el Estado debe demostrar la razonabilidad de la
norma. \\
¿Cómo se resolvió el caso de educación religiosa en Salta? & La Corte, por
mayoría, declaró inconstitucional la asignatura de religión obligatoria en la
educación inicial, por generar discriminación indirecta hacia los alumnos no
católicos. \\
¿Qué rige la negativa a transfusiones de sangre por motivos religiosos? & El
principio de reserva y la libertad religiosa, siempre que sea una decisión
personal, inequívoca y lúcida, no suplible por terceros. \\
¿Qué eje común atraviesa todos los casos de libertad religiosa? & El principio de
reserva del artículo 19 CN. \\
\midrule
\multicolumn{2}{p{0.95\textwidth}}{\textbf{Resumen:} el sostenimiento
institucional del culto católico convive con la exigencia de neutralidad estatal,
reforzada por la doctrina de las categorías sospechosas, que invierte la carga de
la prueba cuando una norma distingue por religión u otros rasgos sensibles. Esta
lógica se proyecta en el fallo sobre educación religiosa en Salta y en los casos
de transfusiones de sangre, educación en el hogar, comunidades menonitas y
objeción religiosa al servicio militar, todos anclados en el principio de reserva
del artículo 19 CN.} \\
\bottomrule
\end{longtable}

\chapter{Libertad de expresión. Prohibición de censura previa. Doctrina Campillay y doctrina de la real malicia. Derecho de réplica}
\label{cap:p14}

\section{La doble dimensión de la libertad de expresión y su importancia democrática}

\begin{defbox}[: Doble dimensión de la libertad de expresión]
Conforme lo sostiene la Corte Interamericana de Derechos Humanos, la libertad de
expresión tiene una \textbf{dimensión individual} —el derecho de cada persona a
expresar y difundir su pensamiento y sus opiniones— y una \textbf{dimensión
social o colectiva} —el derecho de la sociedad a recibir información y opiniones,
entendida como condición para el funcionamiento del sistema democrático y para la
existencia de un verdadero mercado de ideas.
\end{defbox}

Esta doble dimensión explica por qué la libertad de expresión no se protege
únicamente como un derecho individual del emisor, sino como un pilar
estructural del sistema democrático en su conjunto: sin un intercambio libre de
ideas e información, la ciudadanía no puede formar una opinión pública informada,
controlar a sus gobernantes ni participar efectivamente en el proceso político.
Por esa razón, en el marco del Pacto de San José de Costa Rica, el estado de la
libertad de expresión es uno de los elementos con que la Carta Democrática de la
OEA examina si un sistema político efectivamente adhiere a los postulados
democráticos. La conexión con el Capítulo~\ref{cap:p4} es directa: una prensa
libre es, junto con elecciones libres, uno de los canales de participación cuyo
buen funcionamiento el control judicial está llamado a proteger.

\section{La prohibición de censura previa}

\begin{defbox}[: Censura previa]
Todo mecanismo de control estatal que impide o condiciona, \textbf{antes} de que
una expresión se difunda, su publicación o circulación —por ejemplo, exigiendo una
autorización previa, o prohibiendo directamente la exhibición o publicación de un
contenido—. Se distingue de la responsabilidad ulterior, que sanciona \textit{ex
post} las consecuencias jurídicas de una expresión ya difundida.
\end{defbox}

El artículo 14 de la Constitución Nacional reconoce el derecho de ``publicar las
ideas por la prensa sin censura previa'', y el Pacto de San José de Costa Rica
prohíbe expresamente la censura previa en su artículo 13. El sistema
interamericano de derechos humanos es, en este punto, particularmente estricto: el
control estatal sobre la expresión debe ser siempre reparador y posterior, nunca
preventivo, salvo la única excepción admitida —la protección de menores de edad
frente a información sensible o dañina—.

\begin{exbox}
En el caso ``La Última Tentación de Cristo'' (Corte Interamericana de Derechos
Humanos), se condenó a Chile por no asegurar la libertad de expresión y por haber
ejercido censura previa al prohibir la exhibición de esa película en su
territorio, precisamente por tratarse de un control estatal anterior a la
difusión de la obra.
\end{exbox}

\begin{impbox}
Por qué la censura previa es más grave que la responsabilidad ulterior. Un
sistema de responsabilidad ulterior permite que la expresión circule y that la
sociedad la conozca, reservando la sanción para después, si corresponde; la
censura previa, en cambio, impide que la expresión llegue siquiera a existir en el
debate público, lesionando tanto la dimensión individual como la dimensión social
de la libertad de expresión.
\end{impbox}

\section{La doctrina de la real malicia}

\begin{defbox}[: Doctrina de la real malicia (\textit{actual malice})]
De origen norteamericano (\textit{New York Times Co.\ v.\ Sullivan}, Corte
Suprema de EE.\,UU., 1964) y adoptada por la Corte argentina: cuando la
información se refiere a una persona pública, sujeta a un escrutinio más
intenso, el medio de comunicación solo responde por una información inexacta si
quien demanda prueba que el medio conocía de antemano la falsedad de la
información y, aun así, decidió difundirla (dolo, o cuando menos una negligencia
grave en el chequeo de las fuentes). La carga de la prueba recae en quien demanda.
\end{defbox}

Se exige, entonces, un estándar de exigencia mucho mayor que la simple prueba de
la falsedad de lo publicado: debe probarse el conocimiento previo de la no
veracidad y, pese a ello, la decisión de avanzar con la difusión. Con el tiempo,
la doctrina se ha ido aplicando de forma más amplia que a los solos funcionarios
de alto rango, alcanzando también a personas que, sin detentar una función
pública, tienen una notoriedad amplia entre la población.

\section{La doctrina Campillay}

\begin{defbox}[: Doctrina Campillay --- tres presupuestos indistintos]
Desarrollo propio y exclusivo del derecho argentino (fallo ``Campillay'', CSJN,
1986), que prevé tres dispositivos mediante los cuales, con un ejercicio prudente
de la libertad de información, los medios pueden asegurarse la no responsabilidad
por la circulación de información falaz:
\begin{enumerate}[leftmargin=1.4em]
  \item \textbf{Cita de la fuente}: una cita concreta, clara, identificable y de
  buena fe, no una mera alusión vaga.
  \item \textbf{Omisión de la identidad} de la persona sobre la que se informa,
  de modo que no haya daño a su reputación ni a su honor.
  \item \textbf{Modo potencial}: la información se difunde sin asegurar su
  carácter inexcusable o indudable, en tono conjetural.
\end{enumerate}
Con que se concrete fehacientemente uno solo de estos presupuestos, no hay
responsabilidad del medio.
\end{defbox}

\begin{table}[H]
\centering
\caption{Real malicia y doctrina Campillay}
\begin{tabularx}{\textwidth}{@{}lXXX@{}}
\toprule
\textbf{Doctrina} & \textbf{Origen} & \textbf{Presupuesto central} & \textbf{Efecto} \\
\midrule
Real malicia & Jurisprudencia de EE.\,UU., adoptada por la CSJN & Falsedad de la
información y conocimiento previo de esa falsedad por el medio & Habilita la
responsabilidad civil frente a personas públicas \\
Campillay & Jurisprudencia propia de la Corte argentina & Omisión de identidad,
cita de fuente o modo potencial & Exime de responsabilidad al medio que cumple
alguno de estos recaudos \\
\bottomrule
\end{tabularx}
\end{table}

\section{El derecho de réplica: el caso ``Ekmekdjian c/ Sofovich''}

\begin{notebox}[Ampliación externa]
Los apuntes de la cátedra no traían desarrollo del derecho de réplica, tema que el
cronograma de la Comisión 603 ubica en una clase posterior a la fecha de este
primer parcial, pero que la guía de preguntas incluye expresamente. El desarrollo
que sigue se elaboró consultando el texto del fallo ``Ekmekdjian, Miguel A.\ c/
Sofovich, Gerardo y otros'' (CSJN, 07/07/1992, Fallos 315:1492) y comentarios
doctrinarios especializados (véase la bibliografía final).
\end{notebox}

\begin{defbox}[: Derecho de réplica, rectificación o respuesta]
Derecho de toda persona afectada por informaciones inexactas o agraviantes,
emitidas en su perjuicio a través de un medio de comunicación, a efectuar por el
mismo órgano de difusión su rectificación o respuesta, en las condiciones que
establezca la ley (art.\ 14, Convención Americana sobre Derechos Humanos ---
Pacto de San José de Costa Rica).
\end{defbox}

\subsection{Los hechos del caso}

El escritor Dalmiro Sáenz efectuó, en el programa televisivo ``La Noche del
Sábado'', conducido por Gerardo Sofovich, expresiones que Miguel Ángel
Ekmekdjian —abogado y docente de derecho constitucional— consideró gravemente
ofensivas hacia Jesucristo y la Virgen María, y por las que se sintió lesionado
en sus sentimientos religiosos. Ekmekdjian remitió una carta documento al
conductor del programa solicitando que la leyera al aire, en ejercicio del
derecho de réplica; ante la negativa de Sofovich, inició una acción de amparo
fundada en el artículo 33 de la Constitución Nacional y en el artículo 14 del
Pacto de San José de Costa Rica. Tanto la primera instancia como la Cámara
Nacional de Apelaciones en lo Civil rechazaron el amparo, sosteniendo que el
derecho de réplica, aunque reconocido por el Pacto, requería una ley del
Congreso que lo reglamentara para ser operativo —siguiendo el criterio que la
propia Corte había sostenido pocos años antes, en ``Ekmekdjian c/ Neustadt''
(1988)—.

\subsection{La decisión de la Corte Suprema}

La Corte Suprema, revirtiendo su propio precedente anterior, hizo lugar al
amparo y sentó dos reglas centrales:

\begin{enumerate}[leftmargin=1.6em]
  \item \textbf{El derecho de réplica es operativo, sin necesidad de ley
  reglamentaria}. La Corte reinterpretó la frase del artículo 14.1 del Pacto de
  San José —``en las condiciones que establezca la ley''— no como una condición
  suspensiva que subordinara la existencia misma del derecho a que el Congreso
  dictara una ley, sino como una simple referencia a cuestiones de
  implementación práctica (el espacio o el lapso en que debe ejercerse la
  respuesta). El derecho de réplica, sostuvo la Corte, ``existe e integra nuestro
  ordenamiento jurídico sin necesidad de que se dicte ninguna ley''.
  \item \textbf{Primacía de los tratados internacionales sobre el derecho
  interno}. Con apoyo en el artículo 27 de la Convención de Viena sobre el
  Derecho de los Tratados —según el cual un Estado no puede invocar su propio
  derecho interno para justificar el incumplimiento de un tratado—, la Corte
  afirmó la prioridad de las normas de un tratado internacional ratificado por
  el país por sobre la legislación interna. Este criterio, sentado dos años antes
  de la reforma constitucional, anticipó el régimen que consagraría el artículo
  75, inciso 22 CN a partir de 1994.
\end{enumerate}

\begin{impbox}
Requisitos para el ejercicio del derecho de réplica según ``Ekmekdjian''. Para
que proceda, quien reclama debe estar directamente aludido por la información
difundida —no basta con sentirse dolido en términos generales por una expresión
que no lo menciona a él en particular—, y el espacio que ocupe la respuesta debe
ser el adecuado a su finalidad, sin que necesariamente deba coincidir en
extensión y ubicación con la publicación original.
\end{impbox}

\section{Relación entre las cuatro doctrinas de este capítulo}

\begin{table}[H]
\centering
\caption{Cuatro mecanismos frente a la libertad de expresión y sus límites}
\begin{tabularx}{\textwidth}{@{}lXX@{}}
\toprule
\textbf{Mecanismo} & \textbf{Momento en que opera} & \textbf{Función} \\
\midrule
Prohibición de censura previa & Antes de la difusión & Impide todo control
estatal preventivo sobre el contenido a publicar \\
Doctrina Campillay & Al momento de informar & Permite al medio eximirse de
responsabilidad ulterior si cumple ciertos recaudos formales \\
Real malicia & Después de difundida la información, frente a una demanda de una
persona pública & Exige probar el conocimiento previo de la falsedad para que el
medio responda \\
Derecho de réplica & Después de difundida la información, a instancia del
afectado & Permite a la persona aludida responder por el mismo medio, sin
necesidad de acreditar dolo ni falsedad del medio \\
\bottomrule
\end{tabularx}
\end{table}

\begin{impbox}
Una distinción clave: el derecho de réplica no exige demostrar que el medio actuó
con real malicia ni que incumplió los recaudos de Campillay. Es un remedio más
simple y más rápido: basta con estar aludido por una información que se
considera inexacta o agraviante para tener derecho a responder por el mismo
medio, sin necesidad de un juicio de responsabilidad civil previo. Real malicia y
Campillay, en cambio, operan como estándares para determinar si el medio debe
indemnizar el daño causado.
\end{impbox}

\begin{exambox}
La libertad de expresión tiene una doble dimensión —individual y social— que la
convierte en un pilar del sistema democrático, razón por la cual el sistema
interamericano prohíbe de manera casi absoluta la censura previa (art.\ 14 CN y
art.\ 13, Pacto de San José): ningún contenido puede impedirse antes de su
difusión, salvo la protección de menores. El control estatal, en cambio, es
siempre reparador: opera después de la difusión, a través de dos estándares que
matizan la responsabilidad de los medios frente a la persona afectada —la real
malicia, de origen norteamericano, que exige probar que el medio conocía la
falsedad de lo difundido y aun así lo publicó; y la doctrina Campillay, de
creación exclusivamente argentina, que exime de responsabilidad al medio que citó
la fuente, omitió la identidad del aludido o informó en modo potencial—. A esos
dos estándares se suma el derecho de réplica, reconocido como operativo —sin
necesidad de ley reglamentaria— en ``Ekmekdjian c/ Sofovich'' (CSJN, 1992): la
persona directamente aludida por una información inexacta o agraviante puede
exigir responder por el mismo medio, con apoyo directo en el artículo 14 del
Pacto de San José de Costa Rica, sin necesidad de acreditar previamente dolo,
negligencia grave o incumplimiento de recaudo formal alguno por parte del medio.
\end{exambox}

\section{Cuadro Cornell: libertad de expresión}

\begin{longtable}{@{}p{0.38\textwidth}p{0.55\textwidth}@{}}
\toprule
\textbf{Preguntas / palabras clave} & \textbf{Notas / respuestas} \\
\midrule
\endhead
¿Qué es la doble dimensión de la libertad de expresión? & La dimensión individual
(expresar y difundir el propio pensamiento) y la dimensión social o colectiva
(recibir información, como condición de la democracia). \\
¿Qué es la censura previa y por qué está prohibida? & Todo control estatal
anterior a la difusión de una expresión; está prohibida (art.\ 14 CN, art.\ 13
Pacto de San José) porque impide que la expresión llegue siquiera a existir en el
debate público. \\
¿Qué es la doctrina de la real malicia? & La regla de \textit{New York Times v.\
Sullivan}: el medio solo responde ante una persona pública si esta prueba que el
medio conocía de antemano la falsedad de la información y aun así la difundió. \\
¿Cuáles son los tres presupuestos de la doctrina Campillay? & Cita de la fuente,
omisión de la identidad del sujeto, o uso del modo potencial; basta con uno para
eximir de responsabilidad al medio. \\
¿Qué resolvió ``Ekmekdjian c/ Sofovich''? & Que el derecho de réplica del
artículo 14 del Pacto de San José es operativo sin necesidad de ley reglamentaria,
y que los tratados internacionales priman sobre el derecho interno (art.\ 27,
Convención de Viena). \\
¿En qué se diferencia el derecho de réplica de la real malicia y de Campillay? &
No exige probar dolo, negligencia ni incumplimiento de recaudos formales: basta
con estar directamente aludido por una información inexacta o agraviante. \\
\midrule
\multicolumn{2}{p{0.95\textwidth}}{\textbf{Resumen:} la libertad de expresión, en
su doble dimensión, es incompatible con cualquier forma de censura previa, y se
concilia con la responsabilidad de los medios mediante dos estándares
complementarios y posteriores a la difusión: la real malicia, que exige probar el
conocimiento previo de la falsedad cuando se trata de personas públicas, y
Campillay, que exime de responsabilidad con el cumplimiento de recaudos
formales. El derecho de réplica, reconocido como operativo en ``Ekmekdjian c/
Sofovich'' (1992), ofrece a la persona aludida una vía más directa e inmediata de
respuesta, con apoyo en el Pacto de San José de Costa Rica y sin necesidad de
acreditar previamente la responsabilidad civil del medio.} \\
\bottomrule
\end{longtable}

% ============================================================
% SÍNTESIS FINAL Y BIBLIOGRAFÍA
% ============================================================
\chapter*{Síntesis final: cómo se conectan las catorce preguntas}
\addcontentsline{toc}{chapter}{Síntesis final}

La guía de examen no es una lista de temas sueltos, sino un recorrido con una
lógica interna. Puede reconstruirse así:

\begin{enumerate}[leftmargin=1.6em]
  \item Todo empieza con la \textbf{interpretación} (Capítulo~\ref{cap:p1}):
  antes de aplicar la Constitución hay que decidir cómo leerla, y esa elección
  metodológica —textualismo o no textualismo— condiciona cualquier respuesta
  posterior. La distinción entre \textit{holding} y \textit{obiter dictum} es la
  herramienta que permite leer correctamente cualquier precedente citado en el
  resto de la guía.
  \item Una vez que se sabe cómo interpretar, surge la pregunta de \textbf{quién}
  tiene la última palabra sobre esa interpretación cuando una norma parece
  contradecir la Constitución: el \textbf{control de constitucionalidad}
  (Capítulo~\ref{cap:p2}), ejercido en la Argentina de manera \textbf{difusa}
  (Capítulo~\ref{cap:p3}) y con el problema de fondo de que quien lo ejerce —el
  Poder Judicial— no tiene mandato electoral directo, lo que plantea la objeción
  \textbf{contramayoritaria} (Capítulo~\ref{cap:p4}).
  \item Ese control necesita un cauce procesal para llegar hasta la Corte
  Suprema: el \textbf{recurso extraordinario federal}, cuyo núcleo es la
  \textbf{cuestión federal} (Capítulo~\ref{cap:p5}), complementado por la vía
  paralela de la \textbf{arbitrariedad de sentencia} (Capítulo~\ref{cap:p6}) y
  por la facultad de declarar la inconstitucionalidad \textbf{de oficio}
  (Capítulo~\ref{cap:p7}). Sus \textbf{requisitos} formales, comunes y propios
  (Capítulo~\ref{cap:p8}) determinan cuándo procede; si se deniega, queda la
  \textbf{queja} (Capítulo~\ref{cap:p9}); mientras tanto, puede necesitarse una
  \textbf{medida cautelar} (Capítulo~\ref{cap:p10}); y, en supuestos
  excepcionalísimos de gravedad institucional, puede saltearse directamente a la
  Corte mediante el \textit{\textbf{per saltum}} (Capítulo~\ref{cap:p11}).
  \item Finalmente, todo ese aparato de control se pone a prueba frente a los
  \textbf{derechos individuales} más sensibles: el \textbf{principio de reserva}
  y la \textbf{objeción de conciencia} (Capítulo~\ref{cap:p12}), la
  \textbf{libertad religiosa} (Capítulo~\ref{cap:p13}) y la \textbf{libertad de
  expresión} (Capítulo~\ref{cap:p14}), donde vuelven a aparecer, aplicadas a
  casos concretos, las categorías sospechosas (eco de la teoría contramayoritaria)
  y la exigencia de que cualquier restricción estatal sea posterior y nunca
  preventiva (eco de la prohibición de censura previa).
\end{enumerate}

\begin{impbox}
Estrategia de estudio. Al repasar cada pregunta, conviene preguntarse siempre: (i)
¿qué corriente de interpretación está en juego, si la hay?; (ii) ¿qué tipo de
cuestión federal se plantea?; (iii) ¿qué fallo o doctrina la ilustra, y cuál es su
\textit{holding}?; y (iv) ¿cómo se conecta con las demás preguntas de la guía?
Responder mecánicamente cada punto por separado, sin advertir estas conexiones,
es el error más frecuente en un examen de esta materia.
\end{impbox}

\chapter*{Bibliografía y fuentes}
\addcontentsline{toc}{chapter}{Bibliografía y fuentes}

\section*{Material proporcionado por la cátedra}
\addcontentsline{toc}{section}{Material proporcionado por la cátedra}

\begin{itemize}[leftmargin=1.4em]
  \item Apuntes integrados de Derecho Constitucional Profundizado (2026),
  elaborados por Isaac Edgar Camacho Ocampo a partir de la transcripción de las
  clases y los apuntes originales de la materia.
  \item Cronograma de clases de Derecho Constitucional Profundizado, Comisión
  603, 2.\textsuperscript{o} cuatrimestre de 2026 (docentes: Dr.\ Andrés
  Uslenghi y Dr.\ Darío M.\ Luna).
  \item Guía de preguntas del primer parcial, a rendirse el 28 de septiembre de
  2026.
  \item Sola, Juan Vicente, \textit{Manual de Derecho Constitucional}, La Ley,
  2010; \textit{Tratado de Derecho Constitucional}, Lexis Nexis, 2012.
  \item Ely, John Hart, \textit{Democracia y desconfianza}, Universidad de los
  Andes, Bogotá, 1997.
\end{itemize}

\section*{Fuentes externas consultadas para completar puntos no desarrollados en los apuntes}
\addcontentsline{toc}{section}{Fuentes externas consultadas}

\begin{itemize}[leftmargin=1.4em]
  \item Sobre \textit{holding} y \textit{obiter dictum}: Domenech, Ernesto E.,
  ``Dicho sea de paso: obiter dicta'', Universidad Nacional de La Plata
  (SEDICI); ``Distinción entre \textit{holding} y \textit{obiter dicta}'',
  material docente, Universidad Nacional de Mar del Plata; artículo sobre el
  precedente ``Einaudi'' de la CSJN (CONICET).
  \item Sobre el Poder Judicial contramayoritario y la teoría de Ely: Bickel,
  Alexander M., \textit{The Least Dangerous Branch: The Supreme Court at the Bar
  of Politics}, Yale University Press, 1962; Carnota, Walter F., ``Minorías ni
  tan `discretas' ni tan `insulares': la Suprema Corte de los Estados Unidos
  frente a la discriminación hoy'', \textit{Revista Derecho de las Minorías},
  vol.\ 3, 2020; Sola, Juan Vicente, ``Sobre las minorías discretas e insulares y
  \textit{Carolene Products}'', Academia Nacional de Ciencias Morales y
  Políticas; ``¿La Justicia debe ser independiente de la mayoría?'', Chequeado.
  \item Sobre objeción de conciencia: fallo ``Portillo, Alfredo s/ infracción
  art.\ 44 Ley 17.531'' (CSJN, 18/04/1989, Fallos 312:496); fallo ``Bahamondez,
  Marcelo s/ medida cautelar'' (CSJN, 06/04/1993, Fallos 316:479); Godachevich,
  Mariano Gabriel, ``Una aproximación a la objeción de conciencia al servicio
  militar. Aspectos constitucionales (nota al fallo `Portillo, Alfredo',
  CSJN, 18/4/89)'', Facultad de Derecho, UBA; Legarre, Santiago, ``Objeción de
  conciencia sin precedentes'', Universidad Católica Argentina.
  \item Sobre el derecho de réplica: fallo ``Ekmekdjian, Miguel A.\ c/ Sofovich,
  Gerardo y otros'' (CSJN, 07/07/1992, Fallos 315:1492).
\end{itemize}

\begin{notebox}[Laguna no completada]
El cronograma de la cátedra menciona, junto a ``Portillo'' y ``Bahamondez'', un
tercer fallo bajo el apellido ``Sisto'' en materia de objeción de conciencia. No
fue posible identificar y verificar con las fuentes disponibles la carátula
exacta, el tribunal y el contenido de ese precedente. Se recomienda verificarlo
directamente con la cátedra antes de utilizarlo como referencia de examen.
\end{notebox}

\end{document}
